% !TeX program = pdflatex
\documentclass[9pt,lineno]{elife}
% Use the onehalfspacing option for 1.5 line spacing
% Use the doublespacing option for 2.0 line spacing

% Core math and fonts
\usepackage[T1]{fontenc}
\usepackage[utf8]{inputenc}
% \usepackage{lmodern} % avoid resetting the default family when using pdfLaTeX
\usepackage{amsmath}
\usepackage{amsfonts}

% Units and chemistry
\usepackage[version=4]{mhchem} % currently unused (\ce{} not in text; comment out if not needed)
\usepackage{siunitx}           % used extensively: \SI, \SIrange, \si

% Layout and typography
\usepackage{textcomp}  % probably optional; no explicit textcomp symbols seen
\usepackage{microtype}
\setlength{\emergencystretch}{2em}
\hbadness=1500
\vfuzz=1pt
\raggedbottom
\usepackage{geometry}  % likely configured by elife; no explicit geometry commands used
% \usepackage{booktabs}  % keep if you have tables elsewhere

% Figures and color
\usepackage{graphicx}
\usepackage{caption}
\usepackage{subcaption} % not yet used in the snippets; keep if you plan sub-Figures

% Hyperlinks and quotes
\usepackage{hyperref}
\usepackage{csquotes}   % \enquote{} etc. not used yet; optional


\usepackage{todonotes}


\hypersetup{colorlinks=true, allcolors=blue}

\sisetup{detect-all}
\DeclareSIUnit\Molar{M}
\DeclareSIUnit\dpf{dpf}
\DeclareSIUnit\fps{fps}

% \usepackage{todonotes}

%* Article setup
% region Article setup
\title{Establishing Trace Conditioning in Head-Fixed Larval Zebrafish Reveals Predictive Motor Suppression}

\author[1*]{Joaquim A.S.M. Contradanças}
\author[1]{Michael B. Orger}
\affil[1]{Champalimaud Research, Champalimaud Foundation, Lisbon, Portugal.}

\corr{joaquim.contradancas@research.fchampalimaud.org}{JASMC}
% endregion

\begin{document}

% region Abstract
%* Abstract
\maketitle
\begin{abstract}
	Whether larval zebrafish can associate events separated by a stimulus-free interval remains unclear. Trace conditioning, in which a temporal gap separates the conditioned stimulus (CS) and unconditioned stimulus (US), provides a framework for studying how nervous systems link temporally discontinuous events. We developed a head-fixed classical conditioning assay in larval zebrafish that pairs a visual CS with an aversive optovin-mediated US. Larvae acquired both delay conditioning and trace conditioning across a 3-s stimulus-free interval. Learning was expressed as a reduction in tail-movement vigor that emerged during training, was absent in unpaired controls, became restricted to periods preceding the expected US, and extinguished during testing. Using comparable inter-trial intervals, conditioning was not detected at the population level with a 10-s trace interval. These findings establish predictive motor suppression as a behavioral readout of associative learning and provide a foundation for future cellular-resolution studies of neural mechanisms supporting associations across temporal gaps.
	\textbf{Keywords:} head-restrained behavior, larval zebrafish, delay conditioning, trace conditioning, optovin
\end{abstract}

%todo Excluded in this first minimal version: Learner classification; Temporal organization of learning; Temporal estimation / interval timing; 3-hour retention; MK-801; ca8 ablation; Imaging
% endregion

% region Introduction
%* Introduction
\section{Introduction DRAFT}

Learning allows animals to use experience to predict future events and adapt their behavior accordingly. In classical conditioning, an initially neutral conditioned stimulus (CS) acquires behavioral significance through repeated pairing with a biologically salient unconditioned stimulus (US). The resulting conditioned response (CR) provides a quantitative readout of the learned relationship between cue and outcome. Because the timing, identity, and contingency of the stimuli can be experimentally controlled, classical conditioning remains a powerful framework for studying how nervous systems transform sensory experience into predictive behavior (Pavlov, 1927; Rescorla, 1988; Fanselow and Poulos, 2005; Gershman et al., 2010).

The temporal relationship between the CS and US is central to what the animal must learn. In delay conditioning, the CS overlaps with the US, allowing the association to be formed between events that are contiguous in time. In trace conditioning, by contrast, the CS ends before the US begins, leaving a stimulus-free interval that must be bridged behaviorally or neurally for learning to occur. Trace conditioning therefore provides a particularly useful assay for asking whether an animal can associate temporally separated events, rather than merely respond to overlapping stimuli. Establishing such learning in experimentally accessible preparations is important for linking behavioral evidence of temporal association to its underlying circuit mechanisms (Thompson and Steinmetz, 2009; Gilmartin and Helmstetter, 2010; Raybuck and Lattal, 2014; Kryukov, 2012).

Larval zebrafish offer a distinctive opportunity to study associative learning at brain-wide scale. Their small and optically transparent brains make them suitable for cellular-resolution functional imaging across large fractions of the nervous system, while their genetic and pharmacological accessibility enables mechanistic interrogation. At the same time, larvae produce quantifiable motor behaviors that can be monitored with high temporal resolution during controlled stimulation. These features make the model attractive for connecting learning behavior to distributed neural dynamics in an intact vertebrate brain (Ahrens et al., 2013; Portugues et al., 2014; Orger and de Polavieja, 2017; Vanwalleghem et al., 2018; Lovett-Barron, 2021).

Despite these advantages, classical conditioning paradigms in larval zebrafish still face methodological challenges. Existing studies have demonstrated associative learning in this model, but performance can be variable and may depend strongly on developmental stage, preparation, stimulus choice, and behavioral readout. Freely swimming assays preserve natural behavior but are less compatible with stable high-resolution imaging. Fully immobilized preparations improve optical stability but restrict behavioral output and can make learning difficult to measure. Head-fixed preparations, in which the head is stabilized while the tail remains free, offer a compromise between optical access and behavioral readout, but require assays that maintain sufficient baseline movement to detect learned changes in motor output (Aizenberg and Schuman, 2011; Valente et al., 2012; Harmon et al., 2017; Matsuda et al., 2017; Palumbo et al., 2020; Dempsey et al., 2022).

A further gap concerns trace conditioning. Delay conditioning has been demonstrated in larval zebrafish, but whether head-fixed larvae can acquire associations across a stimulus-free interval remains less clear. This question is important because trace conditioning extends the behavioral demands beyond simple temporal overlap between cue and outcome. Demonstrating trace learning in an imaging-compatible larval preparation would provide a foundation for future studies of how temporally separated events are represented and linked across the vertebrate brain. However, such a claim requires a robust assay with appropriate unpaired controls and careful separation of associative learning from changes in arousal, baseline activity, or non-associative responses to repeated stimulation (Aizenberg and Schuman, 2011; Harmon et al., 2017; Dempsey et al., 2022; Raybuck and Lattal, 2014).

Here, we developed a head-fixed larval zebrafish conditioning assay designed to preserve compatibility with future brain-wide imaging while retaining a measurable behavioral output. The assay pairs a visual CS with an aversive optovin-mediated US and quantifies behavior through tail-movement vigor. Optovin provides a temporally controlled stimulus that reliably evokes tail responses when activated by violet light, allowing repeated US delivery without mechanical contact. Because the learned response in this preparation is expressed as a reduction, rather than an increase, in tail movement, the protocol includes a priming and habituation structure intended to maintain an active baseline from which suppressive conditioned responses can be detected (Kokel et al., 2013; Aizenberg and Schuman, 2011; Harmon et al., 2017).

Using this framework, we tested whether larvae acquire both delay and trace associations. We show that head-fixed larvae develop learned suppression of tail-movement vigor after CS onset in delay conditioning and in a short-trace protocol with a 3-s stimulus-free interval. This suppression is absent in matched unpaired controls and diminishes during testing, supporting an associative interpretation. Under comparable conditions, conditioning was not detected at the population level with a 10-s trace interval, suggesting that longer temporal gaps impose a stronger constraint on this assay or preparation. Together, these findings establish predictive motor suppression as a behavioral readout of associative learning in head-fixed larval zebrafish and provide a foundation for future cellular-resolution studies of the neural mechanisms that support temporal association.

% endregion Introduction

% * Results
% region Results

\section{Results}

\subsection{Setup, Stimuli, and Experimental Protocols}
% Intro: motivate optovin as US and present assay design
TRPA1b-mediated photostimulation with optovin can be delivered with precise timing, evoking robust aversive responses in larval zebrafish~\cite{Wee2019}. We therefore hypothesized that optovin stimulation could serve as an effective unconditioned stimulus (US) for associative learning.
% todo confirm if reference is appropriate and if there are other references to add
% Introduce Delay conditioning as initial validation
To test this, we developed a classical-conditioning assay for head-fixed larval zebrafish with precise stimulus control and automatic behavioral tracking. As a validation of the optovin-based US, we first tested larvae under Delay conditioning, a classical-conditioning configuration already shown to elicit robust learning in larval zebrafish ~\cite{aizenbergCerebellarDependentLearningLarval2011}.
% todo add references to previous works on Delay conditioning in zebrafish
% (Aizenberg and Schuman, 2011; Harmon et al., 2017; Matsuda et al., 2017, Palumbo et al., 2020; Zocchi et al., 2025, Dempsey et al., 2022, Hinz et al., 2013; Zocchi et al., 2025)

% region Setup and Stimuli
\paragraph{Setup, CS, and US}
% Describe head-fixed preparation and tail tracking
In the experimental apparatus, a zebrafish larva was head-restrained in an agarose gel, with its eyes and tail free to move. Tail kinematics were recorded by automatically tracking the tail.
% Define CS illumination and duration
White or red LEDs were used as visual stimuli. A pair of LEDs in front of the larvae provided constant basal illumination; this illumination switched off when the pair of LEDs on the right side turned on (represented in green in \FIG{fig1}A) for \SI{10}{\second}, which served as the conditioned stimulus (CS; Materials and Methods).
% Define optovin-based US and behavioral response
Fish medium contained \SI{10}{\micro\Molar} optovin. For conditioning, \SI{100}{\milli\second} violet-light pulses triggered the optovin-mediated stimulation (\FIG{fig1}A; Materials and Methods). This consistently produced vigorous tail responses (Sup. Fig.~1).
% todo add reference to Sup. Fig. 1
% todo cite Wee  al., 2019; Kokel

% region Experimental Protocols
\paragraph{Experimental Protocols}
% Outline the four main phases of the conditioning protocol
Overall, conditioning experiments included four main phases (\FIG{fig1} B; Sup. Fig. 2 A; Materials and Methods):
\begin{enumerate}
  \item Priming/habituation, 4 CS-like stimuli alone to habituate larvae to the CS;
  \item Pre-training, 10 CS alone;
  \item Training, 46 CS-US paired trials (using \SI{100}{\milli\second} long violet light pulses for the US, \textit{long US}) and 4 CS alone (catch trials); and
  \item Testing, 30 CS alone.
\end{enumerate}

% region Novelty and importance of short USs
% Note optovin elevates baseline activity, motivating protocol adjustments
Optovin stimuli not only evoked robust and vigorous tail responses (\FIG{fig1} C,D), but also elevated baseline swimming activity to levels unusual for head-fixed larvae (Sup. Fig.~1).
% todo what are the typical baseline movement rates of head-fixed larvae in the literature? reference supporting this
% todo indicate the numbers of tail movements and time moving
% Explain rationale for adding short USs to prime and sustain activity
This was an important consideration for our behavioral readout, which was based on tail movement vigor, and greatly influenced protocol design (Sup. Fig. 2 A; Materials and Methods). We therefore strategically incorporated \SI{50}{\milli\second} long violet light pulses (\textit{short US}) at key points in the protocol to first prime the larvae into an active behavioral state of high baseline movement rate (during the priming/habituation phase) and then to maintain it throughout the experiment. These short USs were never used for conditioning; they were never paired with the CS across trials and were temporally separated from any CS presentation by >30~s (Materials and Methods).
% todo confirm time between any short US to any CS presentation
% todo swimming activity levels over the experiment (Thesis Figure 2.3 C, D, only for 100 ms, 10 μM)
% Summarize outcome of short-US strategy on baseline stability
Together, these measures maintained a stable, high level of tail activity throughout the experiment (Sup. Fig.~2 B) and, as shown below, were essential for detecting the conditioned response.
% todo (DETAILED PROTOCOL Sup. Fig. 2 A; BASELINE MOVEMENT RATE Sup. Fig. 2 B)







% region Delay Protocols
% Define Delay CS-US timing
Training was characterized by a (long) US presented \SI{1}{\second} before the CS offset (for Delay Conditioning) or some seconds after the CS offset (for Trace Conditioning) (\FIG{fig1} B) in paired CS-US trials during training (\FIG{fig1} XXXX).
% Describe randomized unpaired Control protocol
For the control, we designed a CS-US unpaired condition in which CS presentation timings were randomized during the training phase (\FIG{fig1} B; Materials and Methods) to preclude associative learning. To further minimize the possibility of accidental conditioning, we assigned nearly every Control fish a unique, randomized CS presentation schedule (Materials and Methods).


\subsection{Delay Conditioning with an Optovin-Based US}
% Define learning metric and tail vigor measure
Learning, in this case the capacity to form associations between the visual CS and the optovin-based US, was assessed by measuring changes in tail movement vigor (rate of change of overall tail curvature over time; Materials and Methods) evoked by the CS over the course of conditioning.

% region CRs in Single Fish
\paragraph{Examples of Potential Learners: Qualitative Analysis of the CR}
% Motivate qualitative single-fish inspection for CR emergence
To gain initial insights into learning-related behavioral changes (appearance of conditioned responses, CRs), we first examined data from individual fish qualitatively.
% Describe baseline normalization of trial-wise vigor
For each trial, tail movement vigor was normalized to the corresponding baseline vigor preceding CS onset (\FIG{fig1} D; Materials and Methods). This normalization accounted for natural oscillations in baseline tail movement across fish and trials (\FIG{fig1} D,E; Materials and Methods).
% Explain heatmap visualization for qualitative assessment
The visualization of these vigor change at CS onset values in heatmaps allowed a first qualitative assessment of CS-evoked responses across the experiment.
% Example Delay fish showing learned suppression during conditioning/testing
During the initial pre-train phase, CS presentation did not elicit any detectable modulation of tail vigor. With training, however, a distinct pattern developed in some Delay fish: vigor during the CS--US interval progressively decreased across trials (\FIG{fig1}D). This reduction became evident early in conditioning and remained stable through the initial testing phase, indicating the acquisition and short-term retention of a conditioned suppression response (conditioned response, CR).
% Contrast Control fish lacking consistent CS-locked suppression
In contrast, Control fish did not exhibit a consistent CS-locked modulation of vigor across trials (\FIG{fig1}E). Although randomized US timing in the Control protocol complicates strict CS-aligned analysis during conditioning, the anticipatory suppression observed during early testing in Delay fish was completely absent in Control animals. This absence confirms that the observed suppression reflects associative learning driven specifically by the CS--US contingency, rather than non-associative changes in motor output.

\paragraph{UR}
% Introduce qualitative features of the optovin-evoked UR
Inspection of individual fish data also highlighted features of the unconditioned response (UR) to the optovin-mediated US.
% Update the following comment: % Note post-UR immobility phase and its early-trial prominence
Tail movements immediately following the US during conditioning trials (\FIG{fig1} D--G US-aligned data for control) differed markedly from pre-stimulus activity, showing distinctly higher frequency and amplitude (\FIG{fig1} C), followed by a phase of immobility. This immobility tended to be most pronounced in the initial conditioning trials---likely reflecting the fish's first encounters with the US.
% todo as previously described... (Wee et al., 2019)
% quantify increased freq. and amp.?

% todo heatmap of US aligned vigor with URs to the short USs
% todo show US-aligned data for the individual fish in Fig. 1 in Sup. Fig.

% todo heatmap of raw vigor?


% TODO
\subsection{PLEASE DO NOT CONTINUE. STILL TO BE DONE}

\paragraph{Population-Level Analysis Confirms Learning in Delay Conditioning}

% TODO
\paragraph{Pooled data heatmaps}

To capture these trends, we aligned and aggregated scaled vigor data from multiple fish for each experimental condition relative to CS onset (Materials and Methods).

the CR first appeared near CS offset, close to the expected US time. As conditioning progressed, the CR shifted earlier to occur shortly after CS onset. During testing, the CR was evident in early trials but extinguished over time, leaving only a mild, startle-like response similar to controls.

% IMPROVE
A comparison between the Delay and trace groups and the respective control fish revealed that only the groups that received the CS-US paired during training developed a consistent behavioural change in response to the CS during the training and testing phases (\FIG{fig2} A--C).
During the first 10 training trials, larvae in the Delay condition showed a progressive reduction in movement vigor within the CS--US interval, emerging after several trials and strengthening with continued training. This conditioned response (CR) became increasingly temporally refined: it initially appeared a few seconds after CS onset, then gradually shifted earlier to align more closely with CS presentation. During early test trials, the suppression often extended beyond CS offset, but it gradually extinguished across testing, with vigor returning to baseline as the CS regained a neutral value.
Larvae trained with short temporal gaps---either using an incrementally increasing trace interval (incTrace) or a fixed \SI{3}{\second} trace interval (3sTrace)---displayed robust associative learning comparable to that observed in the Delay condition. In both paradigms, the CR first appeared near CS offset, close to the expected US time. As training progressed, the CR shifted earlier to occur shortly after CS onset. During testing, the CR was evident in early trials but extinguished over time, leaving only a mild, startle-like response similar to controls.
When the trace interval was extended to \SI{10}{\second}, conditioning became substantially more difficult. Heatmaps of pooled scaled vigor revealed only a subtle suppression that developed after several training trials, localized primarily within the trace interval (between CS offset and US onset). This weak and delayed reduction---absent in control fish---suggests that some associative learning can occur even across this extended temporal gap, though with greatly diminished strength.
%! The 10sTrace heatmaps lacked the clear CS-locked suppression seen in shorter traces, showing instead only a modest reduction in vigor during the trace interval.

% todo show heatmaps with percentage of active fish per time bin (Sup. Fig. ???)

% Transition from single-fish examples to population analysis
After qualitatively assessing individual learning trajectories, we next analyzed responses at the population level.

\paragraph{Control}
Across all experimental conditions, the unpaired control groups provided a critical baseline for comparison.
These controls showed no sustained, learning-dependent changes.
CS-aligned scaled vigor in control animals remained relatively constant throughout the entire experiment, apart from a modest but consistent transient increase during CS presentation (visible as a brief intensity rise in the heatmaps). The UR was not apparent in pooled CS-aligned control data because each fish experienced a distinct control protocol, in which US presentations were explicitly unpaired and temporally varied. Consequently, pooling CS-aligned data across control individuals temporally misaligned the URs, causing them to appear diffuse or attenuated in the population heatmaps.
US-aligned scaled vigor remained stable in the period preceding the US, showing no systematic anticipatory modulation---activity (lacked any build-up or suppression leading up to US onset).
% todo Sup. Fig. with US-aligned data for control
Therefore, the training phase in the control group of fish did not induce consistent activity locked to CS onset nor US onset across trials (Fig. ??? and Sup. Fig. ???).






\paragraph{Block-Level Statistics and Extinction}



Our quantitative analysis focused on comparing vigor change at CS onset across three distinct phases, each representing a 5-trial block:
Pre-Train (PT): The final 5 trials of pre-training, establishing baseline responsiveness.
Early Test (ET): The first 5 trials of testing, measuring the expression of the learned CR.
Late Test (LT): The final 5 trials of testing, assessing the extinction of the CR.
% todo Change abbreviations of experimental blocks of trials



% Quantify Delay-group CR expression using predefined 5-trial blocks.
Block-level analysis of the average vigor change at CS onset confirmed robust CR expression in the Delay group. During Early Test (first 5 test trials), vigor significantly decreased at CS onset relative to both the Pre-Train baseline (final 5 pre-train trials) and the Control group. This suppression was transient: by Late Test (final 5 test trials), vigor change returned toward pre-train levels, indicating extinction of the learned response.
% vigor change at CS onset (still need to find a better name for this)

\paragraph{Vigor Change at CS Onset Across Trials}
% Introduce trial-by-trial NV to track the learning trajectory and extinction dynamics.
To further characterize the learning trajectory and extinction dynamics, we quantified vigor change on a trial-by-trial basis.
% Describe trial-by-trial CR emergence, timing refinement, persistence, and extinction.
Trial-by-trial average vigor change at CS onset showed gradual CR acquisition in the Delay group during conditioning. Suppression within the CS--US interval emerged over approximately the first 10 conditioning trials (around overall trial 20). This suppression persisted into the initial testing phase (approximately trials 61--75) and declined across later test trials, consistent with extinction. In contrast, the control group remained near baseline across trials and showed no consistent trial-locked suppression pattern.

% todo Explain choice of CR windows in Materials and Methods (based on Figure 3)

%! todo when saying "rapid" learning, analyze if after a spurious CS-US pairing there is a sudden drop in NV



\paragraph{Vigor Change in US-aligned Data}
% US aligned Control Data
To confirm that the learned response was specifically an anticipatory behavior cued by the CS (and not other environmental stimuli), we analyzed control vigor change at CS onset aligned to the US onset during conditioning.
Control groups for both the Delay and 3sTrace conditions showed vigor change at CS onset fluctuating around 1.0, with no systematic activity patterns emerging before their randomly timed US presentations.
This behavioral stability in US-aligned control data (also exemplified in pooled scaled vigor heatmaps) strongly indicates that only the explicitly presented CS became predictive of the US. This finding ensures that the learned behavioral changes observed in conditioning groups can be confidently attributed to the specific CS-US contingencies.




\paragraph{Control Specificity: Violet Light Without Optovin}
% Specificity control: test whether violet light pulse alone can serve as an effective US.
In a preliminary control experiment, fish experienced Delay-style CS-US pairing with \SI{100}{\milli\second} violet-light pulses but without optovin. Under these conditions, violet light alone did not function as an effective US (Sup. Fig. 3), supporting the specificity of optovin-mediated reinforcement.

\paragraph{Interpretive Bridge: Baseline Stabilization and UR Profile}
% Keep concise interpretation here; detailed mechanistic interpretation can be expanded in Discussion.
The optovin-evoked UR appeared biphasic, with an early high-amplitude motor burst followed by a longer immobility phase. This qualitative profile is compatible with a fast escape-like component followed by a delayed inhibitory state and provides context for why maintaining baseline movement was important for CR quantification.

% todo Sup. Fig. XXXX with UV light only control


% todo Sup. Fig. with US-aligned data for control


% todo show heatmaps with percentage of active fish per time bin (Sup. Fig. XXXX)


% todo Change abbreviations of experimental blocks of trials
% todo ? How much was the median suppression?







%! todo While the main analyses pooled data from fish tested with both white and red CS, a dedicated analysis exclusively on the red CS group confirmed its suitability for subsequent neuroimaging experiments.

% !!!! vigor is undefined in the absence of movement (Materials and Methods). the heatmaps display the average vigor within each time bin.


% !!!! NOTES:

% !!!! 1. Keep the pooled data heatmaps.

% !!!! 2. Try LTM as last panel of Fig 1. Otherwise, move it to a separate Fig.

% !!!! 3. Try Figure of CR dynamics with catch trials with only learners.
% !!!! Also, might consider just pooling all trials of Trace, cropping at US onset.





For the discussion or so
% Clarify analytical relevance of short-US strategy at the protocol stage
This sustained baseline activity was essential for downstream analyses because the conditioned response (CR) in this paradigm is expressed as a suppression of tail movement vigor.







\subsection{Population-Level Analysis Confirms Learning in both Delay and Trace Conditioning}





\subsection{Population-Level Analysis Confirms Learning in 3 s Trace Conditioning}

% Define trace conditioning intervals
In subsequent experiments, we investigated whether larvae could form associations between temporally separated events using trace conditioning protocols. Under these protocols, during Train, CS was temporally separated from US by \SI{3}{\second} (3sTrace) or \SI{10}{\second} (10sTrace; \FIG{fig1} XXXX).


% region Pop. Heatmaps
\paragraph{Population Heatmaps}

% todo 3sTrace

To capture these trends, we aligned and aggregated scaled vigor data from multiple fish for each experimental condition relative to CS onset (Materials and Methods).

the CR first appeared near CS offset, close to the expected US time. As conditioning progressed, the CR shifted earlier to occur shortly after CS onset. During testing, the CR was evident in early trials but extinguished over time, leaving only a mild, startle-like response similar to controls.


When the trace interval was extended to \SI{10}{\second}, conditioning became substantially more difficult. Heatmaps of pooled scaled vigor revealed only a subtle suppression that developed after several conditioning trials, localized primarily within the trace interval (between CS offset and US onset). This weak and delayed reduction---absent in control fish---suggests that some associative learning can occur even across this extended temporal gap, though with greatly diminished strength.
%! The 10sTrace heatmaps lacked the clear CS-locked suppression seen in shorter traces, showing instead only a modest reduction in vigor during the trace interval.

The UR remained evident during training as a bright band of elevated vigor (\textasciitilde\SI{9}{\second} after CS onset in Delay, Sup. Fig. for US-aligned data...).



% Transition from single-fish examples to population analysis
After qualitatively assessing individual learning trajectories, we next analyzed responses at the population level.
To determine whether the response patterns observed in individual fish were consistent across animals, we pooled CS-aligned tail-movement data for each conditioning group.


% !!!! vigor is undefined in the absence of movement (Materials and Methods). the heatmaps display the average vigor within each time bin.
Because vigor was undefined when a fish was not moving, each heatmap shows the mean scaled vigor of the fish that exhibited tail movements within each time bin. The heatmaps therefore describe the intensity of movement among active fish rather than the proportion of fish that moved.

Because vigor is undefined in the absence of movement (Materials and Methods), the heatmaps display the average scaled vigor across fish that exhibited tail movements within each time bin. Consequently, bins in which only a small subset of fish moved represent the average scaled vigor of that active subset, rather than population-wide immobility. Periods when most animals were immobile---such as the post-US immobility phase of the UR following the initial optovin-evoked burst---therefore do not appear uniformly dark, but instead reflect the residual activity of the few individuals that remained moving.

% Delay and 3sTrace
In the Delay and 3sTrace groups (\FIG{fig2} A,B), CS-US pairing produced a progressive reduction in vigor during training. The CR initially emerged near the expected US time and then shifted earlier during conditioning, becoming apparent soon after CS onset. This suppression persisted into the first test trials but weakened during later testing, consistent with extinction.
% todo show heatmaps with percentage of active fish per time bin (Sup. Fig. ???)


% Control
The corresponding control groups showed no sustained CS-locked suppression. Their CS-aligned vigor remained broadly stable across trials, apart from a brief increase during CS presentation. Because US delivery was explicitly unpaired and varied in time across control animals, US-evoked responses were not aligned in the pooled CS-aligned heatmaps.

The unpaired control groups provided a critical baseline for comparison.
These controls showed no sustained, learning-dependent changes.
CS-aligned scaled vigor in control animals remained relatively constant throughout the entire experiment, apart from a modest but consistent transient increase during CS presentation (visible as a brief intensity rise in the heatmaps). The UR was not apparent in pooled CS-aligned control data because each fish experienced a distinct control protocol, in which US presentations were explicitly unpaired and temporally varied. Consequently, pooling CS-aligned data across control individuals temporally misaligned the URs, causing them to appear diffuse or attenuated in the population heatmaps.
US-aligned scaled vigor remained stable in the period preceding the US, showing no systematic anticipatory modulation---activity (lacked any build-up or suppression leading up to US onset). 
% todo Sup. Fig. with US-aligned data for control
Therefore, the training phase in the control group of fish did not induce consistent activity locked to CS onset nor US onset across trials (Fig. ??? and Sup. Fig. ???).

\paragraph{Control}
Across all experimental conditions, the unpaired control groups provided a critical baseline for comparison.


% 10sTrace
The 10sTrace group showed a different pattern. Pooled heatmaps revealed, at most, a weak suppression during the trace interval between CS offset and US onset. This effect was not consistently distinguishable from control responses at the population level, suggesting that the 10-s interval substantially reduced the strength of associative learning.

%! The 10sTrace heatmaps lacked the clear CS-locked suppression seen in shorter traces, showing instead only a modest reduction in vigor during the trace interval.






\paragraph{Block-Level Statistics and Extinction}



Our quantitative analysis focused on comparing vigor change at CS onset across three distinct phases, each representing a 5-trial block:
Pre-Train (PT): The final 5 trials of pre-training, establishing baseline responsiveness.
Early Test (ET): The first 5 trials of testing, measuring the expression of the learned CR.
Late Test (LT): The final 5 trials of testing, assessing the extinction of the CR.
% todo Change abbreviations of experimental blocks of trials



% Quantify Delay-group CR expression using predefined 5-trial blocks.
Block-level analysis of the average vigor change at CS onset confirmed robust CR expression in the Delay group. During Early Test (first 5 test trials), vigor significantly decreased at CS onset relative to both the Pre-Train baseline (final 5 pre-train trials) and the Control group. This suppression was transient: by Late Test (final 5 test trials), vigor change returned toward pre-train levels, indicating extinction of the learned response.
% vigor change at CS onset (still need to find a better name for this)






\paragraph{Vigor Change at CS Onset Across Trials}
% Introduce trial-by-trial NV to track the learning trajectory and extinction dynamics.
To further characterize the learning trajectory and extinction dynamics, we quantified vigor change on a trial-by-trial basis.
% Describe trial-by-trial CR emergence, timing refinement, persistence, and extinction.
Trial-by-trial average vigor change at CS onset showed gradual CR acquisition in the Delay group during conditioning. Suppression within the CS--US interval emerged over approximately the first 10 conditioning trials (around overall trial 20). This suppression persisted into the initial testing phase (approximately trials 61--75) and declined across later test trials, consistent with extinction. In contrast, the control group remained near baseline across trials and showed no consistent trial-locked suppression pattern.

% todo Explain choice of CR windows in Materials and Methods (based on Figure 3)

%! todo when saying "rapid" learning, analyze if after a spurious CS-US pairing there is a sudden drop in NV




% This will be mostly in Sup. Fig.?

\paragraph{Vigor Change in US-aligned Data}
% US aligned Control Data
To confirm that the learned response was specifically an anticipatory behavior cued by the CS (and not other environmental stimuli), we analyzed control vigor change at CS onset aligned to the US onset during conditioning.
Control groups for both the Delay and 3sTrace conditions showed vigor change at CS onset fluctuating around 1.0, with no systematic activity patterns emerging before their randomly timed US presentations.
This behavioral stability in US-aligned control data (also exemplified in pooled scaled vigor heatmaps) strongly indicates that only the explicitly presented CS became predictive of the US. This finding ensures that the learned behavioral changes observed in conditioning groups can be confidently attributed to the specific CS-US contingencies.




\paragraph{Control Specificity: Violet Light Without Optovin}
% Specificity control: test whether violet light pulse alone can serve as an effective US.
In a preliminary control experiment, fish experienced Delay-style CS-US pairing with \SI{100}{\milli\second} violet-light pulses but without optovin. Under these conditions, violet light alone did not function as an effective US (Sup. Fig. 3), supporting the specificity of optovin-mediated reinforcement.

\paragraph{Interpretive Bridge: Baseline Stabilization and UR Profile}
% Keep concise interpretation here; detailed mechanistic interpretation can be expanded in Discussion.
The optovin-evoked UR appeared biphasic, with an early high-amplitude motor burst followed by a longer immobility phase. This qualitative profile is compatible with a fast escape-like component followed by a delayed inhibitory state and provides context for why maintaining baseline movement was important for CR quantification.

% todo Sup. Fig. XXXX with UV light only control


% todo Sup. Fig. with US-aligned data for control


% todo show heatmaps with percentage of active fish per time bin (Sup. Fig. XXXX)


% todo Change abbreviations of experimental blocks of trials
% todo ? How much was the median suppression?







%% endregion Pop. Analysis

STILL NEED TO REARRANGE THE FOLLOWING SECTIONS.


STOP READING




% region Learners vs Non-Learners

% ** Figure 3
% Example Trace fish
Analyzing individual fish allowed a detailed examination of response development, inter-individual variability, and subtle behavioral features that might be obscured when data are aggregated. Among exemplar fish, tail movements varied substantially in timing and duration both within and across individuals, highlighting the importance of population-level analyses for identifying consistent patterns.






Both in 3sTrace and 10sTrace, we found example larvae exhibiting a reduction in vigor following CS presentation, which developed during Train.

In an example 3sTrace fish, in the initial testing trials, the reduction in vigor at CS onset that developed with conditioning persisted well beyond the end of the CS and the expected time of US (\SI{13}{\second} relativeto CS onset).

An example 10sTrace fish displayed a reduction in vigor that was largely confined to the \SI{10}{\second} trace interval and became progressively consolidated with conditioning; in early test trials, this reduction also extended substantially beyond the trace period.


Analyzing individual fish allowed a detailed examination of response development, inter-individual variability, and subtle behavioral features that might be obscured when data are aggregated. Among example fish, tail movements varied substantially in timing and duration both within and across individuals, highlighting the importance of population-level analyses for identifying consistent patterns. After qualitatively assessing individual learning trajectories, we next analyzed responses at the population level.
Analyzing individual fish allowed a detailed examination of response development, inter-individual variability, and subtle behavioral features that might be obscured when data are aggregated.
Among exemplar fish, tail movements varied substantially in timing and duration both within and across individuals, highlighting the importance of population-level analyses for identifying consistent patterns.
The proportion of larvae successfully conditioned in the 3sTrace groups is reduced compared to the Delay group.
Analysis of individual fish trajectories confirmed these population-level trends. Many experimental fish across these three paradigms showed a downward trend in vigor from pre-conditioning through conditioning and into the early test phase, followed by a return towards baseline during late testing. However, noticeable inter-individual variability was present, with some fish exhibiting more pronounced CRs than others.
Noticeable inter-individual variability was present, with some fish exhibiting more pronounced and sustained CRs than others. This variability appeared higher in the 3sTrace condition compared to the Delay group with some 3sTrace individuals showing less distinct learning trajectories.
% todo Sup. Fig. XXXX with individual vigor change at CS onset trajectories

The 10sTrace condition posed the greatest challenge for learning. At the population level, vigor change at CS onset in the experimental group closely resembled that of controls across all phases, indicating a failure to develop a consistent CR.
Individual data confirmed this trend: most 10sTrace experimental fish had vigor change at CS onset trajectories indistinguishable from controls. A small minority (X of XXX fish) displayed a weak tendency toward reduced vigor during late conditioning and early testing, but this effect was subtle and not representative of the group as a whole.
% todo Sup. Fig. XXXX with individual vigor change at CS onset trajectories
% todo count number of learners in 10sTrace group

% endregion Learners vs Non-Learners





%% region CR Dynamics









%!!! talk about this in Figure 4, when characterizing the CR in learners!
A comparison between the Delay and trace groups and the respective control fish revealed that only the groups that received the CS-US paired during training developed a consistent behavioural change in response to the CS during the training and testing phases 
During the first 10 training trials, larvae in the Delay condition showed a progressive reduction in movement vigor within the CS--US interval, emerging after several trials and strengthening with continued training. This conditioned response (CR) became increasingly temporally refined: it initially appeared a few seconds after CS onset, then gradually shifted earlier to align more closely with CS presentation. During early test trials, the suppression often extended beyond CS offset, but it gradually extinguished across testing, with vigor returning to baseline as the CS regained a neutral value.
Larvae trained with short temporal gaps--- fixed \SI{3}{\second} trace interval (3sTrace)---displayed robust associative learning comparable to that observed in the Delay condition. In both paradigms, the CR first appeared near CS offset, close to the expected US time. As training progressed, the CR shifted earlier to occur shortly after CS onset. During testing, the CR was evident in early trials but extinguished over time, leaving only a mild, startle-like response similar to controls.
When the trace interval was extended to \SI{10}{\second}, conditioning became substantially more difficult. Heatmaps of pooled scaled vigor revealed only a subtle suppression that developed after several training trials, localized primarily within the trace interval (between CS offset and US onset). This weak and delayed reduction---absent in control fish---suggests that some associative learning can occur even across this extended temporal gap, though with greatly diminished strength.






\subsection{CR Dynamics Reveal Precise Temporal Expectation of the US in Delay and Short-Trace Conditioning}

To characterize how the CR developed and when it was expressed relative to the expected US, we combined two complementary analyses: (i) block-averaged scaled-vigor traces (10-trial blocks) and (ii) catch-trial traces (CS-only trials embedded in training).
% todo averaged...?

In the delay condition, CRs emerged rapidly and became both stronger and temporally precise with training. Block averages showed an initial suppression of vigor during the latter part of the CS in the earliest training block (train 1), which strengthened and shifted earlier over subsequent blocks. By late training, suppression began shortly after CS onset and peaked before the expected US time. Catch-trial traces confirmed this timing: CR suppression typically began within a few seconds of CS onset, persisted for about \SIrange{1}{2}{\second} after CS offset, and returned to baseline shortly after the expected US time (Fig. X). During early testing (CS alone), the CR remained robust and time-locked to the predicted US; across test blocks the response shortened and weakened, consistent with extinction (Fig. X).

The and 3sTrace groups displayed qualitatively similar dynamics to delay conditioning but with trace-specific structure. When trace intervals were short ( initially used a \SI{0.5}{\second} gap and was progressively extended; 3sTrace used a fixed \SI{3}{\second} gap), block averages showed a suppression that began during the CS and extended into the stimulus-free interval. With training, the suppression intensified and became focused so that pooled catch-trial data peaked within the trace interval at the predicted US time (Fig. X; Suppl. Fig. Y). These patterns indicate that larvae maintained a CS representation across the gap and developed accurate temporal expectation of US onset. As in delay conditioning, CRs in these groups persisted into early testing and then declined.

The 10sTrace group showed clearly weaker and more variable CRs. Block averages indicated only a modest suppression emerging late in training and persisting into the early test blocks; suppression was largely confined to the long trace interval, whereas scaled vigor during the CS itself stayed near baseline (Fig. X). Pooled catch-trial averages revealed a subtle, diffuse reduction in vigor across the extended trace window that often extended beyond the expected US time. Quantitatively, the 10sTrace group exhibited a statistically significant but attenuated reduction in normalized vigor during the early test phase compared with both pre-train baseline and controls (see Methods and Fig. X), consistent with associative learning that is imprecise in time when the CS--US interval is long.

Across conditions, two consistent themes emerged. First, catch trials isolated a temporally precise CR in delay and short-trace paradigms: suppression was anticipatory (appearing before the expected US) and returned to baseline shortly after the expected US time, demonstrating accurate temporal estimation. Second, single-fish responses were heterogeneous---some individuals showed strong, well-timed CRs while others did not---so pooled averages are complemented by exemplar single-trial traces (Suppl. Fig. Y). Finally, CRs diminished across test blocks in all learned groups, reflecting extinction when the CS was repeatedly presented without the US.











% endregion CR Dynamics










%% region NMDAR dependence, ca8 ablation, LTM with puromycin

% region MK-801
% ** Delay and Trace Conditioning Performance with MK-801

\subsection{Associative Learning is Dependent on NMDA Receptor Function}

To probe the molecular basis of associative learning, I examined the role of NMDARs using the non-competitive antagonist MK-801. The experiment comprised two paradigms: delay conditioning, with MK-801 at 0~\si{\micro\Molar}, 10~\si{\micro\Molar}, or 100~\si{\micro\Molar}, and 3sTrace conditioning, with 0~\si{\micro\Molar} and 100~\si{\micro\Molar}. All test solutions contained 0.2\% DMSO. Control groups received explicitly unpaired CS--US presentations to dissociate associative from non-associative effects.

\paragraph{Delay and Trace Conditioning Performance with MK-801}
Blocking NMDARs with high-dose MK-801 (100~\si{\micro\Molar}) completely abolished CR acquisition in both delay and 3sTrace paradigms (Fig. X). In the Early Test phase, normalized vigor in these larvae did not differ from Pre-Train or from unpaired controls, indicating a failure to learn. In contrast, 10~\si{\micro\Molar} MK-801 produced no measurable impairment in delay conditioning (Fig. X), demonstrating a clear dose dependence of NMDAR involvement.
During delay conditioning, larvae exposed to 0~\si{\micro\Molar} or 10~\si{\micro\Molar} MK-801 displayed a consistent reduction in scaled vigor following CS onset from the second training block onward—an expression of a CR that extinguished by the second test block (Fig. X). In Early Test trials, CR duration was slightly prolonged in the 10~\si{\micro\Molar} group (lasting $\approx$~5~s after CS offset) relative to the 0~\si{\micro\Molar} group ($\approx$~2.5~s). By contrast, 100~\si{\micro\Molar} MK-801 disrupted learning entirely: scaled vigor remained at baseline and did not differ from unpaired controls across experimental phases (Figs. X--Y).
In the 3sTrace paradigm, larvae at 0~\si{\micro\Molar} successfully acquired a CR, whereas those treated with 100~\si{\micro\Molar} failed to do so (Fig. X). Early Test normalized-vigor values of the high-dose group were indistinguishable from Pre-Train and control values. This complete loss of associative suppression confirms that NMDAR signaling is required for both contiguous (delay) and temporally discontiguous (trace) learning.
Subtle trends toward reduced vigor in some high-dose delay fish suggested incomplete blockade or partial compensation, but these were inconsistent across individuals. Variability among 100~\si{\micro\Molar} unpaired controls likely reflects direct motor effects of MK-801 rather than genuine learning. Only delay-conditioning larvae were pre-incubated with MK-801 prior to training, a difference that may further enhance sensitivity to the drug in that group.
% todo These results led me to conclude that NMDA receptor function is critical for this form of associative learning in larval zebrafish, as expected from the literature (Hinz et al., 2013).

\paragraph{Motor Effects of MK-801}
Because I measured the CR as a reduction in tail movement vigor, maintaining a high baseline of spontaneous movement was essential for detecting learning. Sustaining sufficient baseline activity after the \SI{3}{\hour} retention interval posed a major challenge. To address this, I introduced a second priming phase immediately before testing, applying the same principle as the initial priming phase in Chapter 3. Based on pilot experiments, I delivered four \SI{100}{\milli\second} US pulses immediately before the test phase to boost tail activity and enable more reliable CR detection.

% endregion MK-801


% region ca8 ablation
% ** Purkinje Cell Ablation Impairs Associative Learning in Delay and Trace Conditioning
\subsection*{Purkinje Cell Ablation Impairs Associative Learning in Delay and Trace Conditioning}

To examine the role of cerebellar Purkinje cells (PCs) in associative learning, we performed chemogenetic ablation using Nifurpirinol in ca8$^+$ larvae. Six groups were compared: unpaired control ca8$^+$ and ca8$^-$; delay ca8$^+$ and ca8$^-$; trace ca8$^+$ and ca8$^-$.

At the population level, no control group and no PC-ablated (ca8$^+$) larvae learned in either the delay or trace paradigms. In contrast, delay ca8$^-$ larvae exhibited robust learning, while trace ca8$^-$ larvae learned to a lesser degree, consistent with the known variability in 3~s trace conditioning performance. These results indicate that intact PCs are required for both delay and trace associative learning, with trace learning showing partial acquisition even in the absence of PCs, reflecting the increased difficulty of the task.

% todo mention increased difficulty of trace conditioning seen in Fig2

% endregion ca8 ablation



% region LTM with puromycin
% ** 3-h Long-Lasting Memory is Not Blocked by Puromycin Under the Current Protocol
\subsection{3-h Long-Lasting Memory is Not Blocked by Puromycin Under the Current Protocol}

To examine the protein synthesis dependence of long-term memory (LTM), I tested delay-conditioned larvae after a 3-h retention interval, with or without the translation inhibitor puromycin (\SI{5}{\milli\gram\per\litre}). Larvae conditioned without puromycin exhibited significant conditioned responses (CRs) during the Early Test phase, showing reduced normalized vigor relative to their Pre-Train baseline and to unpaired controls (Fig.~X). Puromycin-treated larvae displayed a similar reduction in vigor, though this effect did not reach statistical significance, possibly due to the smaller sample size (Fig.~X, top). These results indicate that, under the present experimental conditions, memory expression after 3~h is not blocked by puromycin. This may reflect either a form of long-lasting memory that does not depend on new protein synthesis, or incomplete inhibition of translation at this concentration.

\paragraph{Paradigm Modifications for LTM and Memory Savings Assessment}
To promote LTM formation, we applied spaced training, inserting a \SI{10}{\minute} interval after every 10 training trials (see 5.4. Materials and Methods). A \SI{3}{\hour} stimulus-free retention interval was introduced between training and testing, during which larvae were kept in darkness to minimize sensory stimulation. Baseline tail activity was maintained by a second priming phase immediately before testing, delivering four \SI{100}{\milli\second} US pulses to ensure reliable CR detection.

To test for memory savings, an additional re-training phase (30 paired trials) followed the extinction block (30 CS-only trials). Learning rates during re-training were compared with both the initial acquisition rates and the performance of naive controls trained for the first time.

\paragraph{Three-Hour Memory in Delay Conditioning}
After the 3-h retention interval, delay-trained larvae without puromycin exhibited robust CRs in the Early Test phase, though the response magnitude was smaller than when testing immediately followed training (Fig.~X). Only a subset of individuals (4 of 27) displayed clear CRs at the single-fish level, consistent with partial retention. Puromycin-treated larvae retained comparable CRs, demonstrating that memory at this timescale was not abolished by protein synthesis inhibition. Because dependence on protein synthesis could not be confirmed, I refer to this as \textit{long-lasting memory} rather than canonical LTM.

\paragraph{No Evidence for Memory Savings}
Re-training analysis revealed no significant evidence of memory savings. Previously conditioned larvae reacquired CRs at rates comparable to their initial acquisition and to naive controls trained for the first time (Fig.~X). Naive controls did not reach equally low normalized-vigor values within the first \textasciitilde10 trials, suggesting a possible qualitative difference in early learning dynamics. At the individual level, only a minority of naive controls (4 of 19) developed clear CRs during re-training, indicating that prior extinction did not facilitate faster re-learning under these conditions (Fig.~X).

Together, these results suggest that 3-h post-conditioning memory persists without clear protein synthesis dependence and does not produce detectable savings upon re-exposure to paired stimuli.
% endregion LTM with puromycin

% endregion NMDAR dependence, ca8 ablation, LTM with puromycin





% Imaging figure
%! todo While the main analyses pooled data from fish tested with both white and red CS, a dedicated analysis exclusively on the red CS group confirmed its suitability for subsequent neuroimaging experiments.





% region Figures

% region Figure 1
\begin{figure}
  \begin{fullwidth}

    \includegraphics[width=0.95\linewidth]{Sections/Figures/Fig1/Fig1_test1.png}

    \caption{A refined classical conditioning paradigm reveals robust associative learning in larval zebrafish. Number of fish tested: Delay group, n=36; Control group, n=31.
      % In the heatmaps, colors represent the average/median scaled vigor in each tail movement.
      % (A) Experimental setup. Top-view schematic of the setup during CS presentation. A larva is head-fixed in agarose with eyes and tail free. Three LED pairs illuminate the arena: the right pair delivers the CS; the front pair provides basal illumination when no CS is present. The larva faces away from a larger LED that delivers the US. Diagram is not to scale except for the LEDs around the head, which are shown in their true relative positions. (B) Conditioning protocol. Experimental groups differ only in CS--US temporal contingency during training. In controls, CS presentations occur at random times relative to the US. (C) Example tail angle traces from two training trials (Early and Late) of a representative delay fish. CS onset and offset are marked by green vertical lines; US onset is marked by the purple vertical line. Grey rectangles indicate automatically detected tail movements. (D) Stimulus delivery protocols. Representative stimulus sequences for all experimental conditions. At this resolution, delay and trace protocols appear similar. In the priming/habituation phase, the second and fourth CSs were presented on the larva's left side. Green lines mark CS onset; purple arrows mark US delivery (\SI{50}{\milli\second} for priming/habituation, pre-training, and test; \SI{100}{\milli\second} for training). (E) Example single-fish vigor. Tail movement vigor from an individual undergoing delay conditioning, showing emergence of a conditioned response (CR). (F) Population-level learning. Heatmap of scaled vigor per trial in delay and control groups. Data were binned in \SI{0.5}{\second} intervals. (G) Quantification of learning. Normalised vigor compared across Pre-Train (PT), Early Test (ET), and Late Test (LT) phases, showing retention and extinction. Each dot represents the median vigor across 5 trials. Asterisks indicate significance (Mann-Whitney U and Wilcoxon signed-rank, Holm-BonferronWe corrected). (H) Acquisition curve. Trial-by-trial evolution of normalised vigor. Shaded areas represent \SI{95}{\percent} confidence intervals. (I) Temporal dynamics. Median $\pm$ confidence interval of scaled vigor (\SI{1}{\second} bins), illustrating CR timing and extinction. Green bars mark CS onset and offset. Figure 1. A refined classical conditioning paradigm reveals robust associative learning in larval zebrafish.
      % (A) Experimental setup. Top-view schematic of the setup during CS presentation. A larva is head-fixed in agarose with eyes and tail free. Three LED pairs illuminate the arena: the right pair delivers the CS; the front pair provides basal illumination when no CS is present. The larva faces away from a larger LED that delivers the US. Diagram is not to scale except for the LEDs around the head, which are shown in their true relative positions.
      % (B) Conditioning protocol. Experimental groups differ only in CS--US temporal contingency during training. In controls, CS presentations occur at random times relative to the US.
      % (C) Example tail angle traces from two training trials (Early and Late) of a representative delay fish. CS onset and offset are marked by green vertical lines; US onset is marked by the purple vertical line. Grey rectangles indicate automatically detected tail movements.
      % (D) Stimulus delivery protocols. Representative stimulus sequences for all experimental conditions. At this resolution, delay and trace protocols appear similar. In the priming/habituation phase, the second and fourth CSs were presented on the larva's left side. Green lines mark CS onset; purple arrows mark US delivery (\SI{50}{\milli\second} for priming/habituation, pre-training, and test; \SI{100}{\milli\second} for training).
      % (E) Example single-fish vigor. Tail movement vigor from an individual undergoing delay conditioning, showing emergence of a conditioned response (CR).
      % (F) Population-level learning. Heatmap of scaled vigor per trial in delay and control groups. Data were binned in \SI{0.5}{\second} intervals.
      % (G) Quantification of learning. Normalised vigor compared across Pre-Train (PT), Early Test (ET), and Late Test (LT) phases, showing retention and extinction. Each dot represents the median vigor across 5 trials. Asterisks indicate significance (Mann-Whitney U and Wilcoxon signed-rank, Holm-BonferronWe corrected).
      % (H) Acquisition curve. Trial-by-trial evolution of normalised vigor. Shaded areas represent \SI{95}{\percent} confidence intervals.
      % (I) Temporal dynamics. Median $\pm$ confidence interval of scaled vigor (\SI{1}{\second} bins), illustrating CR timing and extinction. Green bars mark CS onset and offset.
      % Number of fish tested: Delay group, n=36; Control group, n=31.
    }
    \label{fig:fig1}
  \end{fullwidth}
\end{figure}
% endregion Figure 1

% region Figure 2
\begin{figure}
  \begin{fullwidth}
    \includegraphics[width=0.85\linewidth]{Sections/Figures/Fig2/Fig2_test1.png}
    \caption{
      Figure 2. A refined classical conditioning paradigm reveals robust associative learning in larval zebrafish.
      % (A) Experimental setup. Top-view schematic of the setup during CS presentation. A larva is head-fixed in agarose with eyes and tail free. Three LED pairs illuminate the arena: the right pair delivers the CS; the front pair provides basal illumination when no CS is present. The larva faces away from a larger LED that delivers the US. Diagram is not to scale except for the LEDs around the head, which are shown in their true relative positions.
      % (B) Conditioning protocol. Experimental groups differ only in CS--US temporal contingency during training. In controls, CS presentations occur at random times relative to the US.
      % (C) Example tail angle traces from two training trials (Early and Late) of a representative delay fish. CS onset and offset are marked by green vertical lines; US onset is marked by the purple vertical line. Grey rectangles indicate automatically detected tail movements.
      % (D) Stimulus delivery protocols. Representative stimulus sequences for all experimental conditions. At this resolution, delay and trace protocols appear similar. In the priming/habituation phase, the second and fourth CSs were presented on the larva's left side. Green lines mark CS onset; purple arrows mark US delivery (\SI{50}{\milli\second} for priming/habituation, pre-training, and test; \SI{100}{\milli\second} for training).
      % (E) Example single-fish vigour. Tail movement vigour from an individual undergoing delay conditioning, showing emergence of a conditioned response (CR).
      % (F) Population-level learning. Heatmap of scaled vigour per trial in delay and control groups. Data were binned in \SI{0.5}{\second} intervals.
      % (G) Quantification of learning. Normalised vigour compared across Pre-Train (PT), Early Test (ET), and Late Test (LT) phases, showing retention and extinction. Each dot represents an individual fish. Asterisks indicate significance (Mann-Whitney U with Holm-BonferronWe correction).
      % (H) Acquisition curves. Trial-by-trial evolution of normalised vigour. Shaded areas represent \SI{95}{\percent} confidence intervals.
      % (D) Temporal dynamics. Median $\pm$ confidence interval of scaled vigour (\SI{1}{\second} bins), illustrating CR timing and extinction. Green bars mark CS onset and offset.
      % Number of fish tested: Delay group, n=36; Control group, n=31.
      % Figure 2. Associative learning in trace conditioning depends on trace interval duration.
      % (A) Population-level learning. Heatmaps of scaled vigour across trials for incTrace (left), 3sTrace (middle), and 10sTrace (right) groups. Data are binned in \SI{0.5}{\second} intervals. Rows represent individual trials (pre-training at top, test at bottom).
      % (B) Quantification of learning. Normalised vigour compared across Pre-Train (PT), Early Test (ET), and Late Test (LT) phases for each trace paradigm. Each dot represents an individual fish. Asterisks indicate significance (Mann-Whitney U with Holm-BonferronWe correction).
      % (C) Acquisition curves. Trial-by-trial evolution of normalised vigour. Shaded areas represent \SI{95}{\percent} confidence intervals.
      % (D) Temporal dynamics. Median $\pm$ confidence interval of scaled vigour (\SI{1}{\second} bins), illustrating CR timing and extinction. Green bars mark CS onset and offset.
      % Number of fish tested: incTrace, n=34 (test), n=36 (control); 3sTrace, n=79 (test), n=83 (control); 10sTrace, n=22 (test), n=18 (control).
    }
    \label{fig:fig2}
  \end{fullwidth}
\end{figure}
% endregion Figure 2

% region Figure 3
\begin{figure}
  \begin{fullwidth}
    \includegraphics[width=0.75\linewidth]{Sections/Figures/Fig3/Fig3_test1.png}
    \caption{
      Timing thing.
    }
    \label{fig:fig3}
  \end{fullwidth}
\end{figure}
% endregion Figure 3


% region Figure 4
\begin{figure}
  \begin{fullwidth}
    \includegraphics[width=0.65\linewidth]{Sections/Figures/Fig4/Fig4_test1.png}
    \caption{
      Molecular mechanisms of associative learning.
      % (A) NMDA receptor dependence. Effect of MK-801 on delay (top) and trace (bottom) conditioning. Boxplots show the distribution of normalised vigour during Pre-Train, Early Test, and Late Test phases. Each dot represents the median vigour from 5 trials per fish. Boxes span the interquartile range (Q1--Q3), with a dashed line marking the median. Whiskers extend to 1.5 $\times$ IQR beyond Q1 and Q3; outliers beyond this range are plotted individually. Asterisks indicate significance (Mann-Whitney U and Wilcoxon signed-rank, Holm-Bonferroni corrected).
      % Number of fish tested: (A) Delay (top), n(test 0~\si{\micro\Molar})=16, n(ctrl 0~\si{\micro\Molar})=14, n(test 10~\si{\micro\Molar})=15, n(ctrl 10~\si{\micro\Molar})=15, n(test 100~\si{\micro\Molar})=13, n(ctrl 100~\si{\micro\Molar})=13; Trace (bottom), n(test 0~\si{\micro\Molar})=17, n(ctrl 0~\si{\micro\Molar})=13, n(test 100~\si{\micro\Molar})=17, n(ctrl 100~\si{\micro\Molar})=16.
      % (B) Protein synthesis and long-term memory. Delay conditioning in the absence (left) and presence (right) of the protein synthesis inhibitor puromycin (5 mg/L). Vigour suppression during Early Test persists under puromycin, indicating protein-synthesis-independent memory formation. Each dot represents an individual fish. Asterisks indicate significance (Mann-Whitney U and Wilcoxon signed-rank, Holm-Bonferroni corrected).
      % Number of fish tested: (B) left, n(test 0 mg/L)=27, n(ctrl 0 mg/L)=19; (B) right, n(test 5 mg/L)=9, n(ctrl 5 mg/L)=12.
      % (C) Memory savings. Normalized vigour across trials during initial training and retraining after extinction. Shaded areas indicate 95\% confidence intervals, and vertical grey bars mark transitions between pre-training, training, testing, and retraining phases. Left: Delay-trained fish ('Delay Train') compared with control fish undergoing their first CS--US paired training trials after testing ('Control Train'). Right: Delay-trained fish during initial training ('Delay Train') compared with the same fish during retraining after extinction ('Delay Re-Train').
      % Number of fish tested: same as in (B) left.
    }
    \label{fig:fig4}
  \end{fullwidth}
\end{figure}
% endregion Figure 4


% % region Figure 5
% \begin{figure}
% \begin{figure}
% \begin{fullwidth}
% \includegraphics[width=0.65\linewidth]{Sections/Figures/Fig5.png}
% \caption{
% unpaired control ca8% \$\textasciicircum +\$ (N = 6) and ca8\$\textasciicircum-\$ (N = 10); delay ca8\$\textasciicircum +\$ (N = 12) and ca8\$\textasciicircum-\$ (N = 10); trace ca8\$\textasciicircum +\$ (N = 13) and ca8\$\textasciicircum-\$ (N = 9). Exclusion criteria were applied prior to analysis.
% }
% \label{fig:fig5}
% \end{fullwidth}
% \end{figure}
% % endregion Figure 5



% endregion Figures


% region Sup. Figures

% todo Also show heatmap of US aligned data for all (weak and strong) US. And raw vigor

%! todo calculate overall informativeness after Training


% * Sup. Fig. 1

% Optovin reliably elicits rapid, vigorous tail bends upon violet-light stimulation, and these responses are tightly time-locked to the stimulus, resembling escape behaviors triggered by noxious cues such as mustard oil or heat \citep{Kokel2013,Wee2019,weeSocialIsolationModulates2022}. Optovin acts via activation of the TRPA1b channel in zebrafish somatosensory neurons \citep{Kokel2013}, providing a robust and well-defined aversive signal that can be precisely timed and standardized across trials. Based on a combination of literature reports and our own optimization experiments (see below), we selected \SI{10}{\micro\Molar} optovin in \SI{10}{\milli\Molar} Tris-HCl-buffered E3 and \SI{100}{\milli\second} violet-light pulses as the US parameters for the classical conditioning protocols.


% To determine the optimal parameters for using optovin as the US in conditioning experiments, we systematically characterized larval responses to combinations of violet-light pulse durations and optovin concentrations. Each condition was tested in separate groups of \SIrange{8}{9}{\dpf} larvae from different clutches. Protocols consisted of 50 violet-light pulses delivered at \SIrange{1.5}{2.5}{\minute} intervals, matching the planned number and spacing of US presentations during conditioning.

% We first screened conditions motivated by previous studies \citep{Lam2017,Wee2019}, ranging from \SI{10}{\milli\second} pulses without optovin (control) to \SI{50}{} and \SI{100}{\milli\second} pulses in the presence of \SI{10}{\micro\Molar} optovin. For each pulse, tail behavior was quantified in the \SI{5}{\second} interval following pulse onset, using (i) the fraction of fish exhibiting at least one tail movement and (ii) the vigor of those movements as readouts of stimulus aversiveness.

% Only conditions with \SI{10}{\micro\Molar} optovin and pulse durations of \SI{50}{} or \SI{100}{\milli\second} elicited tail movements in \textasciitilde\SI{97}{\percent} of fish across all 50 stimuli. By contrast, \SI{10}{\milli\second} pulses, with or without optovin or with \SI{5}{\micro\Molar} optovin, produced responses in fewer than \SI{75}{\percent} of fish on average and yielded vigor values comparable to those observed in the absence of optovin. These conditions were therefore deemed insufficiently robust for use as a US.

% Between the two effective conditions at \SI{10}{\micro\Molar} optovin, \SI{50}{\milli\second} pulses evoked an average tail-movement vigor of \textasciitilde\SI{21}{\degree\per\milli\second}, whereas \SI{100}{\milli\second} pulses evoked \textasciitilde\SI{25}{\degree\per\milli\second}. In both cases, vigor decreased over the first \textasciitilde25 stimuli and then reached a plateau. We selected \SI{100}{\milli\second} pulses at \SI{10}{\micro\Molar} optovin for conditioning experiments because they produced stronger movements and thus a more salient US.

% We also examined how repeated optovin stimulation affected spontaneous activity. In the 15-s baseline window preceding each pulse, more than \SI{80}{\percent} of larvae exposed to \SI{50}{} or \SI{100}{\milli\second} pulses at \SI{10}{\micro\Molar} optovin displayed at least one tail movement in almost every trial, whereas baseline movement was markedly lower under all \SI{10}{\milli\second} conditions. With longer pulses, the number of movements and the time spent moving increased over repeated exposures, while the vigor of individual bouts remained comparable between the \SI{50}{} and \SI{100}{\milli\second} conditions. These observations indicate that optovin plus violet light not only drives strong US-locked responses but also elevates overall locomotor activity across trials.

% Beyond choosing US duration and optovin concentration, buffering the medium proved essential for consistent stimulus delivery. In preliminary experiments conducted in unbuffered E3, the characteristic yellowish color of optovin faded over several hours, the medium pH rose above 8.0 (relative to the typical larval rearing pH of \textasciitilde7.2 \citep{Westerfield2007,Martins2016}), and behavioral responses to the optovin stimulus weakened over time (data not shown). To stabilize both pH and optovin, we added \SI{10}{\milli\Molar} Tris-HCl to the E3 medium. Qualitative observations indicated that this modification permitted experiments lasting more than \SI{5}{\hour}, with larvae appearing healthier and more active, and with optovin-evoked responses remaining stable (data not shown). Consequently, all subsequent conditioning experiments used \SI{10}{\micro\Molar} optovin in E3 buffered with \SI{10}{\milli\Molar} Tris-HCl, with \SI{100}{\milli\second} violet-light pulses as the US.

% * Sup. Fig. 2

% To assess how paradigm refinements affected baseline activity, I analyzed tail movement during the \SI{15}{\second} pre-CS period across all trials for delay and 3sTrace experiments. Both groups exhibited similar behavioral patterns.
% A high proportion of fish consistently exhibited movement during the pre-CS baseline across all trials (Figures 3.21 A, 3.22 A). On average, over \SI{96}{\percent} of fish in both conditions displayed at least a tail movement during the baseline period. Furthermore, the probability of movement during the pre-CS period remained relatively stable throughout the experiment.
% The average vigor of pre-CS baseline tail movements also remained remarkably stable throughout the experiment for all conditions examined (Figures 3.21 B, 3.22 B), averaging 14--16 deg/ms over the entire experiment.
% Importantly, the refined paradigm led to greater consistency in the overall amount of movement. Both the number of tail movements and the proportion of time spent moving during the pre-CS period were notably stable across all trials. On average, fish displayed a minimum of 9.6 tail movements and spent >\SI{25}{\percent} of the time moving during the baseline across all conditions. These parameters demonstrate reduced trial-to-trial fluctuation compared to the original paradigm (Figure 2.8).
% In conclusion, the refined conditioning paradigm successfully established a consistently high and more stable baseline of tail movement throughout the experiment, underscoring the effectiveness of the priming and habituation phases, along with interspersed US presentations during pre-train and testing phases, in keeping the animals engaged. The enhanced stability, particularly evident in the number of movements and total time spent moving, represents a significant improvement by providing a more reliable behavioral foundation for assessing CRs.
% The results demonstrate that the refined paradigm successfully established and maintained a stable and active behavioral baseline, which is crucial for reliably detecting learning-induced changes in vigor.


% * Sup. Fig. 3
% In this 'violet light only' control, the heatmap appears overall darker (Figure 2.4 B). Although the number of fish was comparable across conditions (30 without optovin and 34 with optovin in the delay conditions), this indicates substantially fewer tail movements overall, resulting in more time bins with no movement from any fish, shown as dark areas in the heatmap. Similarly to the unpaired control group, the scaled vigor heatmap for 'violet light only' group of fish showed no apparent behavioral differences between the pre-CS period and the CS-US interval.


% * Sup. Fig. ???   Validation of priming/habituation phase
% \subsection{Validation of priming/habituation phase}

% To validate the effectiveness of the priming/habituation phase in establishing a high and stable baseline level of locomotor activity, we quantified baseline tail movements in larvae exposed to the delay protocol under four stimulus conditions that differed in the presence or absence of optovin (opt.) and violet light (VL). Larvae (\SIrange{6}{8}{\dpf}) were head restrained and subjected to priming/habituation and pre-training phases identical in structure to those used in the delay conditioning paradigm (14 CSs and 18 USs of \SI{50}{\milli\second} each, with stimuli delivered at \textasciitilde\SI{1}{\minute} intervals and consecutive USs spaced \textasciitilde\SI{2}{\minute} apart). The four groups were: (i) no VL, no optovin; (ii) VL only, no optovin; (iii) optovin only, no VL; and (iv) both VL and optovin.

% Baseline motor activity was quantified during the 15-s pre-CS period of each trial. For each larva and trial, we measured (i) whether at least one tail movement occurred (fraction of larvae with any pre-CS movement), (ii) the number of tail movements, (iii) the fraction of time spent moving, and (iv) the vigor of tail movements, defined as the mean tail-movement vigor across all detected bouts within the pre-CS window. Vigor was only defined during periods of active movement; consequently, variability in vigor estimates (SEM) reflects how many bouts contributed to the mean in each condition and trial.

% Presentation of VL alone, without optovin, did not alter baseline activity relative to the condition lacking both VL and optovin. Optovin alone induced a modest increase in spontaneous swimming. By contrast, the combined presence of optovin and VL produced a pronounced and sustained elevation in baseline activity: in this group, the fraction of larvae exhibiting any pre-CS movement reached \textasciitilde0.88 (compared to <0.60 in all other conditions), and from trials 4--14 larvae performed on average \textasciitilde8.5 tail movements and spent \textasciitilde\SI{22}{\percent} of the pre-CS period moving. Despite these differences in activity levels, the mean vigor of individual bouts was similar across conditions; larger SEM values in groups with lower activity were attributable to the smaller number of movements contributing to each average.

% These results demonstrate that simultaneous application of optovin and violet light robustly increases both the likelihood and extent of larval movement during the priming/habituation and pre-training phases, thereby establishing the desired high and stable behavioral baseline for subsequent assessment of conditioned reductions in tail vigor.



% region Validation of Red CS for Neuroimaging Applications
% ** Validation of Red CS for Neuroimaging Applications
% \paragraph{Validation of Red CS for Neuroimaging Applications}
% While previous analyses pooled data from experiments using both white and red CS, dedicated validation of a red CS was essential for its intended application in neuroimaging studies. A crucial step in early paradigm development, before commencing imaging experiments, was to confirm effective learning with a red CS, thereby addressing potential concerns about differential perception by larval zebrafish compared to white light. A red CS is particularly advantageous for GCaMP-based functional imaging because its wavelength can be more effectively filtered from the green GCaMP emission spectrum, which minimizes signal interference and enhances data quality. The 3sTrace group, with its substantially larger sample size, was chosen for this focused validation. Quantitative analysis of the 3sTrace group exposed exclusively to the red CS provided robust evidence of successful associative learning, followed by complete extinction of the learned response. The successful conditioning with the red CS validates its use for subsequent neuroimaging investigations.
% endregion Validation of Red CS for Neuroimaging Applications



% endregion Sup. Figures


% endregion Results





%* Discussion
% region Discussion

% ** Discussion

In summary, the quantitative analysis of vigor change at CS onset provides strong evidence for successful associative learning and extinction in the Delay, and 3sTrace paradigms. These results, which align with the patterns observed in the scaled vigour plots, confirm that larval zebrafish can robustly form associations with short temporal gaps. However, they also highlight the significant challenge of conditioning over an extended 10-second trace interval, where a clear learning effect was not statistically evident at the population level.

% Learning Rates Comparison
At the population level, larvae in the Delay and 3sTrace paradigms demonstrated a robust and gradual acquisition of the CR. This manifested as a progressive reduction in vigor during the CS presentation as conditioning progressed. This learning was rapid.

The Delay group began to show a statistically significant reduction in vigor change at CS onset relative to controls within the first 10 conditioning trials (approx. overall trial 20).
The 3sTrace group exhibited acquisition within the first 10--15 conditioning trials.
% todo I think it is slower...
This learned suppression persisted into the initial testing phase (approx. trials 61--75) before gradually extinguishing. In stark contrast, the control group showed no CR, with vigor change at CS onset remaining consistent.


% Clarify analytical relevance of short-US strategy at the protocol stage
This sustained baseline activity was essential for downstream analyses because the conditioned response (CR) in this paradigm is expressed as a suppression of tail movement vigor.







\section{Discussion}

We established a robust classical conditioning paradigm for head-fixed larval zebrafish that supports both delay and trace learning and reveals predictive suppression of motor activity as a conditioned response (CR). Using optovin-induced activation of TRPA1b channels as an aversive unconditioned stimulus (US), we achieved reliable associative learning measured through reductions in tail-movement vigor during the conditioned stimulus (CS). This work is the first to demonstrate classical conditioning in larvae using optovin as a US, and the first to establish trace conditioning in this model organism. Together with subsequent pharmacological and memory-retention experiments, these results lay the foundation for dissecting the neural mechanisms of associative learning and temporal prediction using whole-brain imaging.

\subsection{A refined paradigm enables consistent associative learning}

Previous conditioning attempts in restrained larvae were hindered by low learning efficiency, inconsistent control protocols, and unstable baseline behavior. The present paradigm overcomes these limitations using several key refinements. A stationary, temporally uniform CS replaced the earlier dynamic visual cue, isolating temporal learning from spatial confounds. Extending CS duration to ten seconds improved the detectability of gradual vigor changes and aligned the assay with the temporal resolution of calcium imaging. Baseline motor priming and intermittent low-intensity US presentations stabilized spontaneous locomotion, which was essential because the CR expresses as a reduction rather than enhancement of movement.

Catch trials interleaved within training sessions enabled measurement of the full CR time course without contamination by the unconditioned response (UR), and the expanded test phase revealed extinction dynamics. These refinements allowed larvae to develop internally timed responses that became increasingly precise as training progressed. The paradigm also revealed that vigorous optovin stimulation elevates baseline locomotion over several minutes, presumably reflecting a sustained arousal state, which both facilitates detection of vigor suppression and may enhance learning itself.

\subsection{Predictive motor suppression as a conditioned response}

In both delay and trace paradigms, the CR manifested as a graded suppression of tail-movement vigor beginning shortly after CS onset and peaking near the expected US. Although the UR consists of vigorous, escape-like movements, the CR reflects suppression. This apparent mismatch is resolved by recognizing that the UR to optovin is biphasic: an initial reflexive locomotor burst followed by a longer-lasting immobility phase. We interpret the CR as anticipatory recruitment of this immobility component, representing a threat-predictive inhibitory response. Similar multi-component URs have been described across zebrafish paradigms, where CR form often mirrors the dominant UR phase expressed under the relevant behavioral constraints. Under head restraint, where escape is impossible, immobility becomes the adaptive defensive mode. Freely swimming larvae may instead acquire a locomotor-facilitation CR, which should be tested in future work.

CR expression exhibited substantial inter-individual variability. Some larvae showed strong vigor suppression; others showed none. This could reflect insufficient baseline movement, alternative CR modalities (e.g., eye or heart-rate changes), or true non-learning. Expanded kinematic and physiological metrics will help reveal hidden forms of learning and reduce apparent non-learner rates.

\subsection{Temporal precision and the limits of trace conditioning}

The progressive shift of vigor suppression toward the expected US demonstrates that larvae acquire both the CS--US contingency and its temporal structure. Delay conditioning produced the strongest and most temporally precise CRs. Trace conditioning with short gaps (3~s) or gradually increasing gaps also supported learning, though with extended suppression into the trace interval. At longer intervals (10~s), responses were weak or non-significant, revealing a boundary on temporal bridging.

Selecting the trace interval required balancing two opposing factors: longer intervals probe the limits of working memory but reduce learning effectiveness, whereas shorter intervals promote robust learning but reveal little about temporal computation. A fixed 6~s trace interval may serve as an optimal compromise for imaging experiments, extending beyond 3~s yet avoiding the failures observed at 10~s, while matching the slower kinetics of GCaMP indicators.

Interval structure also influences informativeness. In the 10~s trace design, the interval between the US and the next CS contracts, reducing predictive value. Future paradigms should equate US--CS intervals across trace lengths to ensure fair comparisons. Training protocols that shape the trace interval across sessions---progressively extending it into the 10~s range---may reveal whether learning can be strengthened by gradually stretching the temporal demand.

\subsection{Comparison with established conditioning models}

This zebrafish paradigm integrates features of eyeblink conditioning (EBC) and lick-suppression conditioning (LSC). Like EBC, the CR is a graded, precisely timed modulation of motor output; like LSC, it appears as suppression relative to an elevated baseline. Maintaining high baseline locomotion was essential, mirroring the requirement for sustained licking to measure suppression in LSC.
Despite differences in effector systems, the learning trajectories---acquisition, timing refinement, and extinction---suggest that conserved computational principles govern predictive inhibition across vertebrates.

\subsection{NMDAR dependence and pharmacological constraints}

Pharmacological manipulations revealed that conditioning in both paradigms depends on NMDAR signaling. Delay conditioning persisted at 10~\si{\micro\Molar} MK-801 but was disrupted at 100~\si{\micro\Molar}. Trace conditioning failed entirely at 100~\si{\micro\Molar}, indicating greater vulnerability to NMDAR blockade, consistent with mammalian work showing that trace learning requires sustained CS representations mediated by NMDAR-dependent persistent activity.

However, high-dose MK-801 also produced severe motor abnormalities: altered baseline locomotion, disrupted UR structure, and loss of forward swims. These effects complicate interpretation because systemic drug administration collapses cognitive and motor contributions. Nonetheless, the qualitative differences between delay and trace groups at the same dose argue that the impairments are not purely motoric. A rigorous dose--response curve, including trace conditioning at 10~\si{\micro\Molar}, is needed to determine whether the difference is categorical or graded. More precise methods---local drug delivery, inducible genetic knockdown, or targeted optogenetic/chemogenetic perturbations---will be essential for isolating NMDAR-dependent circuit mechanisms.

\subsection{Neural substrates and implications}

The emergence of trace conditioning implies that larvae maintain CS information across a temporal gap, a computation attributed in mammals to hippocampal, prefrontal, and cerebellar loops. In zebrafish, the pallium and cerebellum are strong candidates for analogous roles. Preliminary two-photon imaging indicates that a subset of putatively cerebellar neurons becomes CS-responsive only after learning. Ongoing and planned experiments using nitroreductase-mediated ablation, laser lesions, and temporally precise optogenetics will test the causal contributions of these populations to delay and trace learning. Whole-brain imaging during training will further reveal how predictive associations are represented and transformed across the vertebrate brain.

\subsection{Memory retention, protein synthesis, and the nature of the 3-h memory}

Using the refined delay protocol, CRs remained detectable after a 3-h retention interval. However, this memory was not abolished by 5~mg/L puromycin. Spaced training conferred no advantage over massed training, consistent with a memory dominated by short-term mechanisms rather than protein-synthesis-dependent long-term memory (LTM). Several explanations are possible: the 3-h memory may not require protein synthesis; drug penetration or timing may have been insufficient; or the memory relies on atypical mechanisms.

Preliminary data at 10~mg/L puromycin suggest reduced CR expression at 3~h, but the experiment was underpowered and lacked unpaired controls. A systematic dose--time study with molecular verification of translational blockade will be required for a decisive classification. Extending retention intervals beyond 3~h remains logistically challenging due to stress, fasting, and fatigue during prolonged restraint. A two-day protocol involving overnight release proved unreliable, underscoring the need for new designs that permit feeding and recovery between sessions.

\subsection{Reacquisition and the absence of memory savings}

Reacquisition after extinction did not proceed faster than initial learning, indicating no behavioral evidence of memory savings. This could reflect true absence of savings, profound extinction, insufficient behavioral sensitivity, or an STM-dominated regime at \SI{3}{\hour}. Interestingly, na\"ive controls exposed only to paired training during re-training showed lower CR magnitudes, suggesting that prior unpaired CS--US exposure may produce \enquote{learned irrelevance}, which suppresses subsequent learning. Future memory-savings designs should manipulate re-training parameters, include explicit learned-irrelevance controls, and exploit neural readouts to detect latent memory traces.

\subsection{Additional limitations and future directions}

A key limitation was statistical power, particularly for 10~s trace conditioning and LTM assays. Larger cohorts will allow more precise estimation of effect sizes, individual differences, and subtle learning dynamics.

Further work is needed to disentangle temporal expectation from CS offset in delay conditioning by systematically decoupling CS duration from US timing. Modulating CS length or shifting US delivery within the CS will reveal whether CR termination is tied to interval timing or driven by CS offset.

Repeated optovin stimulation elevates baseline locomotion, potentially reflecting arousal or sensitization. US-only sessions and variations in inter-US intervals will clarify whether arousal, TRPA1 kinetics, or photochemical effects drive this baseline shift. Non-aversive priming stimuli, such as optomotor responses or whole-field luminance changes, may provide cleaner baselines.

Future paradigms could explore discriminative conditioning and operant learning using the same hardware by assigning different CS LEDs to different contingencies. Computational modeling of learning curves, extinction trajectories, and CR timing will provide a principled framework for inferring the algorithms governing associative learning.

\subsection{Concluding remarks}

This work establishes a versatile conditioning platform in larval \textit{Danio rerio} compatible with whole-brain imaging. Larvae learned temporally structured associations, expressed predictive motor suppression, bridged trace intervals up to 3~s, and exhibited NMDAR-dependent learning. Memory persisted for at least 3~h but could not yet be classified as protein-synthesis-dependent LTM, and memory savings were undetected. The combination of optical accessibility, genetic tractability, and compact brain size makes larval zebrafish uniquely suited for mapping how associative memories are formed, maintained, and expressed at the level of entire neural circuits. Future work integrating behavior, pharmacology, whole-brain activity mapping, and targeted perturbations promises to reveal the circuit motifs and computational strategies that underlie associative learning and temporal prediction.


% endregion


% \input{Sections/materialsandmethods.tex}

%* Methods and Materials
% region Methods and Materials

% ** Materials & Methods

\section{Materials and methods}

Because vigor is undefined in the absence of movement (Materials and Methods), the heatmaps display the average scaled vigor across fish that exhibited tail movements within each time bin. Consequently, bins in which only a small subset of fish moved represent the average scaled vigor of that active subset, rather than population-wide immobility. Periods when most animals were immobile---such as the post-US immobility phase of the UR following the initial optovin-evoked burst---therefore do not appear uniformly dark, but instead reflect the residual activity of the few individuals that remained moving.
% ! in the heatmaps, x values indicate the value of the right side of the bin
The only difference between the protocols of the conditions we tested was the relative timing of the CS and US during this training phase.
Two groups of larvae experienced trace conditioning, where in CS--US paired trials during the training phase CS and US were separated in time with no overlap by \SI{3}{\second} or \SI{10}{\second} trace intervals, respectively.
To investigate the impact of trace interval duration, which is known to be a critical factor, we conditioned different groups of larvae under trace conditioning with trace intervals of \SI{3}{\second} (3sTrace) and \SI{10}{\second} (10sTrace).
% (Gallistel and Gibbon, 2000; Raybuck and Lattal, 2014)





\paragraph{Summary of Sample Sizes}
For reference, the main experimental groups and sample sizes are summarized in Table~\ref{tab:groupsizes}.

\begin{table}[ht]
  \centering
  \caption{Number of fish per group for the refined conditioning paradigm.}
  \label{tab:groupsizes}
  \begin{tabular}{lcc}
    \toprule
    Condition & Test (paired) & Control (unpaired) \\
    \midrule
    Delay     & 36            & 31                 \\
     & 34            & 36                 \\
    3sTrace   & 79            & 83                 \\
    10sTrace  & 11            & 15                 \\
    \bottomrule
  \end{tabular}
\end{table}




\subsection{Animals and husbandry}

All experiments were performed on larval zebrafish (\textit{Danio rerio}) maintained under standard laboratory conditions. For behavioral conditioning assays without concurrent brain imaging, we used nacre (mitfa\textsuperscript{\(-/-\)}) larvae and, in some experiments, nacre larvae expressing pan-neuronal GCaMP (Tg(elavl3:GFF); Tg(10xUAS:GCaMP6fEF05)). For experiments involving chemogenetic ablation of Purkinje cells, nacre larvae carrying Tg(ca8:epNtr-tagRFP) were used, with Ntr\textsuperscript{\(+\)} and Ntr\textsuperscript{\(-\)} siblings tested in parallel.
% TODO clearly define the genotypes used in each experiment.
% TODO In Sup. Fig. with original paradigm, where the no-optovin control is mentioned, state that the larvae were of the Tübingen (TU) strain.
% TODO For the original no-optovin control experiment: Larvae were raised exclusively in E3 medium, which was exchanged daily after hatching. To provide nutritional supplementation and promote activity, a few milliliters of live rotifers (Brachionus plicatilis, cultured as in \citet{Martins2016}) were added daily after the medium change. This medium was replaced daily and the stock was stored for up to four days in an open container to ensure adequate oxygenation.
% TODO Add this information for the original experiment in the Sup. Fig. with the no-optovin control.

Unless otherwise indicated, larvae were obtained from natural pairwise or group matings of adult breeders. Embryos were collected in the morning and kept in groups of 80--100 per 150~mm Petri dish containing E3 embryo medium (\SI{5}{\milli\Molar} NaCl, \SI{0.17}{\milli\Molar} KCl, \SI{0.33}{\milli\Molar} CaCl\textsubscript{2}, \SI{0.33}{\milli\Molar} MgSO\textsubscript{4}; pH adjusted to 7.2 with HCl) at \SI{28}{\degreeCelsius} under a 14/10~h light--dark cycle starting at 08:00, following standard protocols \citep{Westerfield2007,Martins2016}. The sex of larvae is not morphologically distinguishable at the ages used (\SIrange{6}{9}{\dpf}). % chktex 8
% TODO In some experiments (those requiring screening for fluorescence-positive larvae), the density was not controlled.

From \SI{0}{\dpf} to \SI{3}{\dpf}, larvae were maintained in standard E3 medium. From \SI{3}{\dpf} onward, larvae were reared in ``green water''. Green water was prepared by diluting one volume of a stock co-culture of rotifers (\textit{Brachionus plicatilis}, cultured as in \citet{Martins2016}; \textasciitilde900~rotifers/mL) and microalgae (e.g. \textit{Chlorella} sp.) in four volumes of fresh E3. This medium was replaced daily, and the stock was stored for up to four days in an open container at room temperature to ensure adequate oxygenation.

% TODO Add this information for the delay and trace experiments done in the initial experiments.

For pharmacological experiments with the NMDA receptor antagonist MK-801, larvae had the genetic background mitfa\textsuperscript{\(-/-\)}; Tg(elavl3:GFF); Tg(10xUAS:\hspace{0pt}GCaMP6fEF05). In the delay-conditioning experiment, larvae were \SIrange{5}{7}{\dpf} and tested in E3/Tris-HCl containing \SI{0.2}{\percent} DMSO (solvent for optovin and MK-801) and MK-801 at final concentrations of 0, 10, or \SI{100}{\micro\Molar}. Prior to training, fish were incubated for \SI{30}{\minute} in their respective test solutions under constant illumination from the frontal LEDs. In the trace-conditioning experiment, larvae were \SIrange{6}{7}{\dpf}, had the same genetic background, and were tested in E3/Tris-HCl containing 0 or \SI{100}{\micro\Molar} MK-801 without prior incubation.

% TODO Add this information for the ca8 ablation experiment.
% TODO ca8 fish were incubated at 4~\dpf overnight in small Petri dishes, without controlling larval density.

For long-term memory experiments with the protein synthesis inhibitor puromycin, subjects were \SIrange{5}{7}{\dpf} larvae. Their genetic background was mitfa\textsuperscript{\(-/-\)}; Tg(elavl3:GFF); and Tg(10xUAS:GCaMP6fEF05). Fish were tested either in control conditions or in the presence of \SI{5}{\milli\gram\per\liter} puromycin, which was added at the start of the experiment without pre-incubation and maintained until the end of the retention interval \citep{hinzProteinSynthesisDependentAssociative2013}.

% TODO Add this information for the imaging experiments.

All animal husbandry and experimental procedures conformed to the European Union Directive 2010/63/EU on the protection of animals used for scientific purposes. Protocols were reviewed and approved by the Champalimaud Foundation Ethics Committee and the Portuguese national authority for animal welfare, Direção-Geral de Alimentação e Veterinária (DGAV; permit no.~019774, in accordance with former Directive 86/609/EEC).


\subsection{Animal preparation}

For experiments, larvae were either head restrained on the day of testing or embedded the day before and maintained overnight. Larvae were immobilized in \SIrange{1.5}{2}{\percent} (w/v) low-melting-point agarose (Fisher Bioreagents) dissolved in E3 medium
buffered with \SI{10}{\milli\Molar} Tris-HCl (E3/Tris-HCl), following established protocols \citep{Portugues2011}. Embedding was performed in \SI{35}{\milli\meter} Petri dishes whose base was coated with SYLGARD\texttrademark~184 elastomer (Dow Corning). After the agarose solidified, the dish was filled with E3/Tris-HCl medium. Agarose surrounding the tail and eyes was carefully removed with a scalpel to permit unrestricted tail and eye movement. Stainless steel insect pins (0.20~mm, Austerlitz 01.03.04) were used to secure the agarose to the SYLGARD-coated dish base, preventing drift. Immediately before the experiment started, the E3/Tris-HCl medium was replaced with E3/Tris-HCl containing \SI{10}{\micro\Molar} optovin (Alomone Labs) and \SI{0.1}{\percent} (v/v) dimethyl sulfoxide (DMSO; Sigma-Aldrich), and larvae were positioned on the acrylic platform of the behavioral setup, with the body axis aligned to the LED array and the tail centered in the camera field of view. All experiments were performed during the light phase of the 14/10~h light--dark cycle.
% todo say "as shown in Fig. 1 A"


\subsection{Behavioral setup}

All experiments were conducted in four identical, custom-built behavioral setups housed inside a light-proof enclosure to prevent external light contamination. Each setup consisted of a transparent acrylic platform on which the embedded larva was positioned and an array of 3 pairs of light-emitting diodes (LEDs) mounted \textasciitilde\SI{6}{\centi\meter} surrounded the fish.
% todo say "as shown in Fig. 1 A"

Two setups used white LEDs (Multicomp Pro, MP000435) and two used red LEDs emitting at \SI{660}{\nano\meter} (Kingbright, L-1334SRT). The white LEDs were driven with a \SI{51}{\ohm} series resistor to ground, and the red LEDs with a \SI{200}{\ohm} series resistor. The luminous intensity of the CS LEDs was not further calibrated.

For delivery of the US, each setup was equipped with a high-power 1-A violet LED emitting at \SIrange{405}{410}{\nano\meter} (LED Engin, LZ1-10UB00-01U8). These LEDs were driven by dedicated constant-current LED drivers (LEDdynamics, 3023-D-E-1000) gated by control signals from the microcontrollers. Output power was calibrated with a photodiode power meter to \textasciitilde\SI{40}{\milli\watt} measured at \SI{2.5}{\centi\meter} from the LED surface, and pulse-width modulation was used to ensure stable light intensity across trials.

Larval behavior was monitored from above with an infrared-sensitive high-speed camera (Allied Vision Mako U-029B) fitted with a \SI{16}{\milli\meter} fixed-focal-length lens (Edmund Optics 59-870) and an \SI{830}{\nano\meter} long-pass filter (Thorlabs FGL830) to block visible light from the CS and US LEDs. The fish were illuminated from below by an array of \SI{850}{\nano\meter} infrared LEDs (Würth Elektronik 15400385A3590), with the IR light passing through diffuser paper to provide homogeneous background illumination for tail tracking.

All four setups were controlled from a single host computer via dedicated external microcontroller boards (Arduino Uno R3). In each rig, three 8-bit shift registers (e.g. SN74HC595) provided independent control over the 24 CS LEDs. The timing and gating of the violet US LEDs were likewise controlled by the Arduinos through the LED driver inputs. Custom C\# software (Sardine framework) handled communication with the microcontrollers, acquisition of camera frames, and real-time logging of stimulus events. The correct operation of all LEDs and cameras was verified daily before experiments.

% todo show scheme with the exact dimensions of the setup


\subsection{Conditioned and unconditioned stimuli}

% todo measure the luminous power of the CS LEDs?
% todo confirm references of the LEDs used
% todo confirm the wavelength range
% todo Scheme of setups
% todo	Relative positions of all LEDs
% todo	Angle and distance of the UV LED
% todoBlack tube used to calibrate UV



\subsubsection{Conditioned stimulus}

In all experiments, the conditioned stimulus (CS) was delivered using three pairs of LEDs: one pair in front of the larva and one pair on each side. The CS consisted of a \SI{10}{\second} illumination of the right-side LED pair. During the habituation phase, two CS presentations were instead delivered from the left side to habituate larvae to CS-like stimuli.

Outside CS periods, the front pair provided constant basal illumination. During a CS, only a single lateral pair was active at any given time and the front pair was switched off, minimizing abrupt changes in luminance at CS onset and offset. Exactly two LEDs were illuminated at a time, ensuring that the fish were never in complete darkness. Maintaining continuous illumination avoids dark-induced reductions in arousal or sleep-like states \citep{Prober2006} and prevents the CS from being confounded with abrupt light--dark transitions that can elicit large-angle startle turns \citep{Burgess2007}.

Two LED colors were used across experiments. In two setups, the CS was generated with white LEDs (400--800~nm), which fall within the optimal spectral-sensitivity range of larval zebrafish photoreceptors and thus provide a salient visual stimulus \citep{Hughes1998}. In the remaining setups, red LEDs emitting at \SI{660}{\nano\meter} were used. Red light lies largely outside the main sensitivity range of larval zebrafish, making it perceptually weaker, but it is readily separable from the green fluorescence of GCaMP (\textasciitilde\SI{510}{\nano\meter}) in imaging experiments, thereby minimizing optical interference. We therefore tested the conditioning paradigm with both white and red CSs: white light to maximize behavioral salience in non-imaging experiments, and red light to evaluate the feasibility of using a spectrally separated CS compatible with whole-brain GCaMP imaging.



\subsubsection{Unconditioned stimulus}

The unconditioned stimulus (US) consisted of a brief violet-light pulse delivered in the presence of optovin. In all conditioning experiments, larvae were immersed in E3 medium buffered with \SI{10}{\milli\Molar} Tris-HCl (E3/Tris-HCl) containing \SI{10}{\micro\Molar} optovin (Alomone Labs) and \SI{0.1}{\percent} (v/v) dimethyl sulfoxide (DMSO; Sigma-Aldrich). The US was generated by a high-power violet LED emitting at \SIrange{405}{410}{\nano\meter} (LED Engin, LZ1-10UB00-01U8) positioned in front of the larva and driven by a constant-current driver as described above. Except during the optimization experiments described below, US presentations during conditioning consisted of \SI{100}{\milli\second} violet-light pulses (\textit{long US}) or \SI{50}{\milli\second} pulses (\textit{short US}) during the priming, pre-training, and testing phases.

\paragraph{Optovin stimulus optimization}

In pilot experiments, we qualitatively compared unbuffered E3 medium with E3 buffered with \SI{10}{\milli\Molar} Tris-HCl. For each condition, we monitored medium pH and the visible stability of optovin over the duration of an experimental session (more than \SI{3}{\hour}), and we assessed whether larvae remained suitable for experiments lasting several hours. Based on these tests, we selected \SI{10}{\micro\Molar} optovin in E3 buffered with \SI{10}{\milli\Molar} Tris-HCl as the standard medium for conditioning experiments, combined with violet-light pulses of either \SI{50}{\milli\second} or \SI{100}{\milli\second} as the US.
% April 2021
% Optovin degrades and changes pH.
% At the beginning of the day, 10 μM optovin had a pH of 7.47, measured ~1 h after preparation.
% More transparent solution is optovin solution that went through trace conditioning (1000-ms trace period; first round of trace experiments; fixed duration of US) and then one block of delay with a different fish had a pH of 8.22 by the end of the day.
% (The smaller tube was used just to measure the pH.)
% The more yellowish solution is optovin prepared with the other one but that was always kept away from light and never in contact with fish.
% At the end of the day its pH was around 8.68.
% The responses of the fish in these experiments designed to test the degradation of optovin with time (?) showed that some kind of degradation might be really happening, explaining the increasingly worse behavior throughout the long experiments. Instead of degradion of optovin, the too-high pH might explain why fish dont respond reliably.
% On the 3rd and 4th May 2021, did some trace conditioning experiments using a 500-ms trace period and optovin solution buffered with 10 mM Tris-HCl (concentration before adding optovin). The fish seemed to be considerably healthier throughout the ~6-h experiment. They were very active throughout and responded almost 100\% of the times to optovin, only missing in the last block very few stimuli.


To establish suitable parameters for using optovin as the US, we systematically varied violet-light pulse duration, optovin concentration, and medium composition. Each condition was tested in separate groups of \SIrange{8}{9}{\dpf} larvae drawn from different clutches. Protocols consisted of 50 violet-light pulses delivered at \SIrange{1.5}{2.5}{\minute} intervals, matching the planned number and spacing of US presentations during conditioning.

We tested violet-light pulse durations of \SI{10}{\milli\second}, \SI{50}{\milli\second}, and \SI{100}{\milli\second} in the presence of 0, \SI{5}{\micro\Molar}, or \SI{10}{\micro\Molar} optovin, as well as vehicle-only controls. For each pulse, tail behavior was quantified in a post-stimulus window by measuring the occurrence and intensity of tail movements. In addition, we quantified spontaneous tail movements in a pre-stimulus baseline window to assess how repeated optovin stimulation affected ongoing locomotor activity.


\subsection{Classical conditioning protocols}

% todo Show schemes of all protocols
% todo Protocol Makers
% todo Histograms showing that the multiple controls provided really randomly unpaired CS and US

\subsubsection{General structure}

% todo see Sup. Fig. 2 A
All conditioning experiments followed a structured sequence of phases designed to establish baseline behavior, induce associative learning, and assess CRs. Over multiple iterations of experiments, we developed a general protocol structure that optimized behavioral performance and accommodated imaging-compatible genotypes. The experiment was divided into four phases: priming/habituation, pre-training, training, and testing.

-----
% This was moved from the Results to here
% Describe priming/habituation structure and short-US placement
To promote an active behavioral state similar to that achieved during conditioning, the experimental protocols started with a priming phase of 15 short US presentations, delivered before the pre-train phase. The habituation to the CS happened during the priming by interleaving CS-like stimuli with the short USs. Short USs were also interspersed throughout both the pre-train and testing phases to prevent a reduction in the baseline movement rate during these periods of the experiment critical for assessing conditioning.
% todo (DETAILED PROTOCOL Sup. Fig. 2 A)
% todo (Thesis Figure 2.8)
-----

The priming/habituation and pre-training phases together comprised 14 CSs and 18 USs of \SI{50}{\milli\second} each, delivered at \textasciitilde\SI{1}{\minute} intervals, with consecutive USs spaced \textasciitilde\SI{2}{\minute} apart. The initial 15 USs interleaved with 4 CSs (including CSs from the left side) constituted the priming/habituation phase, which was designed to elevate and stabilize baseline locomotor activity and to habituate larvae to CS-like stimuli.

The subsequent 10 CSs and 3 USs formed the pre-training phase, in which a \SI{50}{\milli\second} US was delivered every 2--3 CSs to further maintain an active behavioral state while measuring baseline CS responses.

The training phase consisted of 50 CS presentations, of which 46 were paired with a \SI{100}{\milli\second} US (delay or trace conditioning) and 4 were CS-only catch trials (see below). Consecutive CSs and consecutive USs were spaced \textasciitilde\SI{2}{\minute} apart.
During the training phase of all delay and trace conditioning protocols, larvae experienced 50 CS presentations. 46 of these CSs were paired with a \SI{100}{\milli\second} US and 4 were CS-only catch trials. The catch trials occurred on CS trials 11, 25, 39, and 45. The first catch trial was therefore presented only after at least 10 CS--US pairings, based on the observation that CRs emerged within this window in pilot delay-conditioning experiments. Catch trials were interspersed with paired CS--US trials rather than presented consecutively, minimizing extinction of the CR. Within each batch of protocols, catch trials were scheduled at identical times relative to the start of the experiment across matched delay, trace, and control protocols, allowing direct comparison of behavior under different CS--US contingencies.

The testing phase mirrored the pre-training structure but was extended to 30 CS presentations interleaved with 14 \SI{50}{\milli\second} USs (one US every 2--3 CSs), allowing CR expression and extinction to be monitored over a larger number of CSs.

After the last CS of the testing phase, a final \SI{500}{\milli\second} US was delivered as a health and viability check.

\subsubsection{Training}
% Delay and trace conditioning
During the training phase of all conditioning protocols, larvae experienced 50 CS presentations. In delay and trace conditioning, 46 of these CSs were paired with a \SI{100}{\milli\second} US, and 4 were CS-only catch trials (see below).

In the delay-conditioning paradigm, the US onset occurred \SI{9}{\second} after CS onset, lasting \SI{100}{\milli\second}.

Three trace protocols were tested: (i) an incremental-trace protocol () in which the trace interval was \SI{0.5}{\second} for the first 10 paired training trials, \SI{3}{\second} for the last 10 paired trials, and increased linearly from \SI{0.5}{\second} to \SI{3}{\second} across the intermediate trials; (ii) a constant-\SI{3}{\second} trace protocol (3sTrace), in which the interval between CS offset and US onset was fixed at \SI{3}{\second} for all paired trials; and (iii) a constant-\SI{10}{\second} trace protocol (10sTrace), in which the US occurred \SI{10}{\second} after CS offset.

Delay conditioning was performed in \SIrange{6}{9}{\dpf} larvae, incremental-trace conditioning in \SIrange{6}{8}{\dpf} larvae, and constant-\SI{3}{\second} and constant-\SI{10}{\second} trace conditioning in \SI{6}{\dpf} and \SI{7}{\dpf} larvae, allowing comparison with established larval fear- and aversive-conditioning paradigms \citep{aizenbergCerebellarDependentLearningLarval2011,matsudaGranuleCellsControl2017}.

\subsubsection{Control protocols}

Unpaired control protocols were designed to match the overall number and duration of CS and US presentations in the corresponding delay and trace conditioning protocols while eliminating any systematic temporal contingency between CS and US. To achieve this, the absolute timing of all US presentations relative to the start of the experiment was kept identical to that in the matched delay or trace protocol. By contrast, the timing of CS presentations during the training phase was randomized independently for each control protocol.

Control protocols were generated in batches using a custom Python script. For each batch, a delay protocol was first created; the US timings from the 46 CS--US paired trials in the training phase were then reused to construct the corresponding trace and control protocols. In control protocols, the 46 CS presentations in the training phase were redistributed over the same time window with the following constraints: (i) a minimum interval of \SI{5}{\second} between consecutive CS onsets and (ii) a minimum interval of \SI{0.5}{\second} between any CS onset and any US onset. CS and US could  overlap by chance, but no fixed temporal relationship was imposed. Nearly every control fish was tested with a unique randomized CS schedule, and subsequent analyses confirmed that neither individual control protocols nor their combination exhibited any learnable temporal structure between CS and US.

US times were not themselves randomized. Randomizing US timings (alone or together with CS timings) would have led to clusters of USs separated by less than \SI{1.5}{\minute}, which was avoided because the optovin US is strongly aversive and repeated closely spaced presentations can markedly alter locomotor state. By maintaining the same US schedule as in the delay and trace protocols while randomizing CS onsets only, control conditions matched overall US exposure without introducing excessive US clustering.

Subsequent analysis confirmed that these control protocols, whether assessed individually or as a group, lacked any consistent temporal structure between CS and US that could support associative learning.
% todo show these plots

\subsubsection{Priming/habituation and pre-training}

A priming/habituation phase was introduced before pre-training to establish a high and stable baseline level of locomotor activity and to habituate larvae to CS-like visual stimuli. Together, the priming/habituation and pre-training phases comprised 14 CSs and 18 USs of \SI{50}{\milli\second} each, delivered at \textasciitilde\SI{1}{\minute} intervals, with consecutive USs spaced \textasciitilde\SI{2}{\minute} apart.

The priming/habituation phase consisted of the first 15 USs interleaved with 4 CSs. During this phase, CSs were presented both from the right (training CS) and from the left, exposing larvae to CS-like stimuli that did not consistently predict the US and thereby reducing the likelihood of startle responses or premature associative learning. The subsequent pre-training phase comprised 10 CSs and 3 \SI{50}{\milli\second} USs, with one US delivered every 2--3 CSs. Inter-stimulus intervals in pre-training were shortened to \textasciitilde\SI{1}{\minute} compared to the \textasciitilde\SI{3}{\minute} ITIs of the original paradigm, reducing total experiment duration while maintaining an active behavioral state.

The testing phase mirrored the structure of pre-training but was extended to 30 CSs interleaved with 14 \SI{50}{\milli\second} USs, again with one US presented every 2--3 CSs. This design allowed the full time course of CR expression and extinction to be measured across a larger number of CS trials. After the final CS of the testing phase, a \SI{500}{\milli\second} US was delivered as a health and viability check; only larvae that exhibited at least one tail movement within \SI{5}{\second} of this final US were retained for analysis (see Inclusion criteria).

\subsubsection{Long-term memory and puromycin}

Long-term memory (LTM) experiments with puromycin used \SIrange{5}{7}{\dpf} larvae of the genetic background mitfa\textsuperscript{\(-/-\)}; Tg(elavl3:GFF); Tg(10xUAS:GCaMP6fEF05). Fish were tested either in control conditions or in the presence of \SI{5}{\milli\gram\per\liter} puromycin, which was added at the start of the experiment without pre-incubation and maintained until the end of the retention interval. The delay-conditioning paradigm described above was modified in several respects. First, the initial training phase was converted from a massed to a spaced configuration: after every 10 training trials, a \SI{10}{\minute} break was inserted in addition to the standard \textasciitilde\SI{2}{\minute} inter-trial interval, and catch trials were introduced during training at trials 16, 25, 36, and 45. Training and testing were separated by a minimum \SI{3}{\hour} retention interval, during which larvae were kept in the dark and undisturbed. Immediately before LTM testing, the medium was exchanged for fresh \SI{10}{\micro\Molar} optovin in E3/Tris-HCl without puromycin, while visual stimulation was minimized. Prior to LTM testing, a priming phase was introduced, consisting of four \SI{100}{\milli\second} US presentations spaced \textasciitilde\SI{2}{\minute} apart, beginning \SI{10}{\minute} after the front LEDs were turned on. The testing phase followed the delay-conditioning procedure described above. After testing, larvae underwent a re-training phase of 30 trials identical to the original (non-spaced) delay paradigm, with two catch trials at trials 11 and 25. In addition to the general inclusion criteria defined below, only larvae that exhibited at least one tail movement within \SI{5}{\second} of every US during the re-training phase were included in the final analysis. Normalized vigor was calculated using a CR interval of \SIrange{0}{9}{\second} after CS onset. In the 0~\si{\milli\gram\per\liter} puromycin condition, 19 control fish and 27 delay fish were analyzed; in the \SI{5}{\milli\gram\per\liter} condition, 12 control fish and 9 delay fish were analyzed. Based on the predefined inclusion criteria, 2 control and 6 delay fish were excluded from the 0~\si{\milli\gram\per\liter} groups, and 1 control and 4 delay fish were excluded from the \SI{5}{\milli\gram\per\liter} groups.

\subsection{Validation of priming/habituation phase}

To validate the effectiveness of the priming/habituation phase in establishing a high and stable baseline level of locomotor activity, we quantified baseline tail movements in larvae exposed to the delay protocol under four stimulus conditions that differed in the presence or absence of optovin (opt.) and violet light (VL). Larvae (\SIrange{6}{8}{\dpf}) were head restrained and subjected to priming/habituation and pre-training phases identical in structure to those used in the delay conditioning paradigm (14 CSs and 18 USs of \SI{50}{\milli\second} each, with stimuli delivered at \textasciitilde\SI{1}{\minute} intervals and consecutive USs spaced \textasciitilde\SI{2}{\minute} apart). The four groups were: (i) no VL, no optovin; (ii) VL only, no optovin; (iii) optovin only, no VL; and (iv) both VL and optovin.

\subsection{Chemogenetic ablation of Purkinje cells}

To assess the contribution of Purkinje cells (PCs) to associative learning, we performed targeted ablation of cerebellar PCs using the nitroreductase (Ntr)--Nifurpirinol (NfP) chemogenetic approach. Ntr catalyzes the reduction of the prodrug NfP into a cytotoxic compound, resulting in selective cell death in Ntr-expressing populations. We used the Tg(ca8:epNtr-tagRFP) line, in which enhanced Ntr (epNtr) fused to tagRFP is expressed under the PC-specific ca8 enhancer, enabling selective ablation of PCs following NfP treatment. The transgenic construct was generated by injecting the ca8:epNtr-tagRFP plasmid together with tol2 mRNA into one-cell-stage TL embryos. Founders showing RFP expression restricted to the cerebellar Purkinje cell layer were outcrossed to TL fish to establish a stable line.

For ablation experiments, progeny from Tg(ca8:epNtr-tagRFP)\textsuperscript{+/-} outcrosses were screened for RFP expression in the cerebellum at \SI{5}{\dpf}. Ten RFP-positive (Ntr\textsuperscript{+}) and ten RFP-negative (Ntr\textsuperscript{-}) larvae were kept together in the same Petri dish to ensure identical handling. At \SI{6}{\dpf}, most of the water was replaced with \SI{30}{\milli\liter} of \SI{5}{\micro\Molar} NfP solution in fish water containing \SI{0.2}{\percent} DMSO (diluted 1:500 from a \SI{2.5}{\milli\Molar} stock solution). Petri dishes were incubated overnight (\SI{16}{\hour}) at \SI{28}{\degreeCelsius} in complete darkness to prevent NfP photodegradation. All handling steps, including NfP preparation and fish transfer, were performed under low-light conditions. After incubation, larvae were rinsed thoroughly three times in fresh fish water to remove residual NfP and DMSO, then allowed to recover for \SI{24}{\hour} before behavioral testing at \SI{7}{\dpf}. After the experiments, larvae were re-screened for RFP to confirm group identity. RFP-negative siblings served as treatment controls. These experiments complement previous work implicating cerebellar Purkinje cells and cerebellar output neurons in larval zebrafish motor learning \citep{harmonDistinctResponsesPurkinje2017,najacSynapticVarianceAction2023}.

\subsection{Data acquisition and online tracking}

Behavioral data acquisition, stimulus presentation, and online tail tracking were controlled by custom-written software in C\# (Microsoft) built on the Sardine experimental framework (Orger lab; A. Laborde and A.\,L. Martins, manuscript in preparation) and using the OpenCV image-processing library. The software communicated with the Arduino microcontroller boards that controlled the CS and US LEDs, acquired image frames from the high-speed camera in each setup, and logged all stimulus and behavioral data streams to disk in real time.

Tail kinematics were tracked online during video acquisition at \SI{700}{\fps}. For each frame, the software identified the coordinates of 15 equidistant points along the tail, from the midpoint of the swim bladder to the tail tip, using a modified version of a previously described tracking algorithm \citep{Marques2018}. The angles between consecutive tail segments defined by these points were computed in real time and stored together with the raw tracking coordinates. This provided a continuous, multi-segment description of tail posture suitable for subsequent computation of tail-movement vigor.

Communication with the Arduino Uno R3 microcontroller boards was handled by a dedicated Sardine module. The start and end of each stimulus (CS and US) were triggered by the microcontrollers and reported back to the host computer over a serial connection. Upon reception of these messages, the C\# software recorded the event identity and a millisecond-resolution system timestamp, enabling subsequent alignment of stimulus events with the camera and tail-tracking data streams.

\subsection{Behavioral data preprocessing and synchronization}

Behavioral data were preprocessed offline using custom Python (version~3.13) scripts based on standard scientific libraries. The primary goal of this pipeline was to temporally synchronize three data streams---camera acquisition logs, tail-tracking data, and stimulus event logs---and to provide a clean, annotated dataset for subsequent analysis.

\paragraph{Camera data integration and synchronization}

The camera log file, which contained the system timestamp corresponding to the arrival of each image frame at the host computer, was used to reconstruct precise frame-acquisition timings and to identify any dropped frames. The inter-frame interval (IFI) was initially estimated as the median time difference between consecutive frame log entries. Periods of stable recording were then identified at the beginning and end of each experiment as contiguous sequences of IFIs matching the estimated framerate within a tolerance of $\pm 0.005$~ms. The first frame in the initial stable region was designated as the temporal reference frame. The mean IFI computed between the first and last stable regions was taken as the operational framerate of the camera (expected to be \textasciitilde\SI{700}{\fps}). Frame reception latencies were analyzed to detect any dropped frames; experiments in which at least one frame was missing were excluded from further analysis.

\paragraph{Tail-tracking data integration and temporal refinement}

Raw tail-tracking data, consisting of the time series of 15 tail-segment angles, were synchronized with the camera frames using shared frame IDs present in both the image acquisition and tracking logs. To ensure temporal consistency across experiments, the tail-angle data were interpolated to a standardized sampling rate of \SI{700}{\fps} using linear interpolation and the operational framerate estimated from the camera log. All data recorded prior to the designated reference frame (corresponding to the initial seconds of the experiment) were discarded. Frame acquisition timestamps were then recalculated relative to the reference frame and the operational framerate, yielding an estimate of the actual camera acquisition time for each behavioral data point. The fixed but unknown delay between image capture at the camera sensor and frame reception by the host computer was not explicitly corrected.

\paragraph{Stimulus protocol integration}

Precise timing information for all CS and US presentations was extracted from the stimulus log file, in which millisecond-resolution timestamps were recorded by the host computer upon receiving trigger messages from the Arduino microcontrollers. These stimulus-event timestamps were aligned with the synchronized behavioral data streams using the common system time base, allowing each CS and US to be mapped onto the corresponding points in the tail-angle and camera time series. As with the camera data, the inherent communication delay between stimulus onset at the hardware level and logging by the host computer was an unknown constant and was therefore not corrected explicitly.

\subsection{Tail movement detection and vigor}

Following synchronization of the behavioral and stimulus data streams, tail kinematics were further processed to detect discrete tail movements and to quantify their vigor.

For each frame, the cumulative sum of the 15 tail-segment angles was computed, providing a scalar measure of overall tail curvature. To reduce high-frequency noise, the cumulative angle signal was first smoothed spatially using a rolling average across tail segments with a 3-segment window. Temporal smoothing was then applied using a rolling average with a \SI{10}{\milli\second} window across time (equivalent to 7 frames at \SI{700}{\fps}). This combination of spatial and temporal filtering produced a cleaner curvature trace suitable for subsequent vigor and bout detection.

Tail-movement vigor was defined as the sum of angular speeds across all 15 tail segments. Angular speed for each segment was computed as the first temporal derivative of its cumulative angle time series. The absolute angular speeds of the 15 segments were then summed for each frame, yielding a single scalar ``tail-movement vigor'' value that captured the intensity of movement in any part of the tail.

Discrete tail movements (bouts) were segmented from the continuous vigor time series using a modified version of a previously established algorithm \citep{Marques2018}. A baseline-corrected activity measure was first computed by subtracting, at each frame, a rolling minimum (20-frame centered window) from a rolling maximum (400-frame centered window) of the vigor signal. This operation enhanced transient peaks relative to slower baseline fluctuations. Frames in which this baseline-corrected vigor exceeded a primary threshold of \SI{4}{\degree\per\milli\second} were initially flagged as candidate movement frames. Candidate segments separated by gaps of fewer than 10 frames were merged into single bouts. Any resulting bout with a total duration shorter than 40 frames was discarded as too brief to constitute a tail movement. For each remaining candidate bout, the peak angular speed was computed; bouts whose peak vigor did not reach a secondary threshold of \SI{1}{\degree\per\milli\second} were excluded, ensuring that only sufficiently strong and sustained events were classified as tail movements.

Vigor values for all time points falling outside the detected bouts were set to NaN, reflecting the fact that vigor was defined only during periods of active movement. For all experiments except the initial US-response characterization, the vigor time series was further smoothed using a non-centered 10-frame rolling average (with \texttt{min\_periods=1} to retain early-bout samples) and then downsampled by retaining every 10th frame, reducing the effective sampling rate to \SI{70}{\hertz}.

For subsequent analysis, each experimental trial was defined as the \SI{90}{\second} window centered on CS onset (from \SI{-45}{\second} to \SI{+45}{\second} relative to CS onset). The vigor data were segmented into these trials and grouped according to experimental phase (pre-training, training, testing) and conditioning protocol (delay, trace variants, unpaired control).

\subsection{Derived behavioral measures}

A series of derived measures was computed from the segmented vigor traces to quantify conditioned changes in tail movement.

\paragraph{Scaled vigor (CS- and US-aligned)}

To compare movement intensity across animals, tail-movement vigor was scaled on a per-trial basis relative to a 15-s baseline. For CS-aligned analyses, the baseline was defined as the \SI{15}{\second} period preceding CS onset; for US-aligned analyses, it was the \SI{15}{\second} period preceding US onset. For each trial, the median vigor in the baseline window was subtracted from the vigor at each time point, and the resulting difference was divided by the same baseline median. This produced a dimensionless ``scaled vigor'' trace that reflected proportional changes in movement intensity relative to baseline. Trials in which no tail movement occurred during the entire baseline window had undefined (NaN) baseline medians and were excluded from analyses involving scaled vigor. For each experimental condition, scaled-vigor traces from individual fish were aggregated by computing the median across fish at each time point.



\paragraph{Baseline measure..and responses to UV..}

Baseline motor activity was quantified during the 15-s pre-CS period of each trial. For each larva and trial, we measured (i) whether at least one tail movement occurred (fraction of larvae with any pre-CS movement), (ii) the number of tail movements, (iii) the fraction of time spent moving, and (iv) the vigor of tail movements, defined as the mean tail-movement vigor across all detected bouts within the pre-CS window. Vigor was only defined during periods of active movement; consequently, variability in vigor estimates reflects how many bouts contributed to the mean in each condition and trial.


\paragraph{Heatmaps and visualization}

For visualization of pooled behavior, aggregated scaled-vigor traces were often binned in either \SI{0.5}{\second} or \SI{1}{\second} intervals by averaging all samples within each bin. For CS-aligned heatmaps of pooled data, the binned scaled vigor was min--max normalized within the pre-CS baseline window using the 10th and 90th percentiles of baseline values as the minimum and maximum, respectively, and then clipped to the [0, 1] interval. For individual-fish heatmaps, tail-movement vigor within each bout was represented as a step function with constant value equal to the bout-averaged vigor. These time series were min--max scaled using the 20th and 80th percentiles of the \SI{40}{\second} pre-CS baseline as minimum and maximum and clipped to [0, 1].

To visualize changes in scaled vigor over time across blocks of trials, binned scaled-vigor traces were baseline-corrected by subtracting the median scaled vigor in the 15-s pre-CS window for each block or catch trial. For experimental blocks, the resulting traces were clipped to ranges such as [$-2$, 2]; for pooled catch-trial data, narrower ranges such as [$-1.5$, 1.5] were used. For plotting scaled vigor per individual catch trial, per-fish binned traces were again min--max normalized within the 15-s pre-CS window (10th/90th percentiles as min/max), median-centered, and clipped to [$-1.5$, 1.5]. Scaled vigor was undefined in trials with only one or no baseline movements.

\paragraph{Normalized vigor and CR (CR) intervals}

To summarize responses per trial, a ``normalized vigor'' metric was computed as the ratio of mean raw vigor in a paradigm-specific CR interval to the mean raw vigor in the 15-s pre-CS baseline. Because vigor was defined only during active movement, all NaN samples (immobility) were excluded from these means, so normalized vigor captured changes in movement intensity conditional on movement occurring.

CR intervals were defined a priori based on scaled-vigor time courses and qualitative heatmaps, with the goal of capturing anticipatory suppression while avoiding contamination by the US-evoked response. For CS-aligned analyses, CR intervals were defined as follows: for Delay, \SIrange{0}{9}{\second} after CS onset (covering the CS while excluding the UR); for, \SIrange{0}{10.5}{\second} after CS onset (covering the entire CS and the initial \SI{0.5}{\second} trace interval, matching the early training trials); for 3sTrace, \SIrange{0}{13}{\second} after CS onset (covering the CS and the \SI{3}{\second} trace interval); and for 10sTrace, the \SI{10}{\second} trace interval between CS offset and expected US onset. For US-aligned analyses, normalized vigor was computed using the same procedure but aligning trials to US onset and defining CR intervals as negative time windows preceding the US (e.g. \SIrange{-9}{0}{\second} for Delay and \SIrange{-13}{0}{\second} for 3sTrace).

\paragraph{Pooled vigor}

For some analyses, pooled vigor traces were used instead of per-trial normalized values. In these cases, for each trial and for each block of interest (e.g.\ late pre-training, early test, late test), vigor traces were scaled to their own pre-CS baseline (typically by subtracting and/or dividing by the baseline median), and then averaged or median-combined across trials and fish within the same condition. This approach provided smooth time-resolved estimates of how tail-movement intensity evolved relative to baseline under different conditioning protocols.

\subsection{Inclusion criteria}

Because tail-movement vigor was defined only during periods of active movement (with vigor set to NaN during immobility), normalized and scaled-vigor measures could only be computed for trials in which larvae moved in both the baseline and CR (CR) intervals. To ensure that group-level analyses reflected reliably engaged animals rather than sporadically active individuals, we applied predefined behavioral inclusion criteria before pooling data across fish. Only larvae that consistently responded to the US and generated sufficient tail movements across key trial blocks were retained.

Unless otherwise indicated (e.g.\ in the priming/habituation validation experiment), a fish was included in the main conditioning analyses only if all of the following criteria were met:
\begin{itemize}
	\item It exhibited at least one tail movement within \SI{5}{\second} following every US in the training phase.
	\item It exhibited at least one tail movement within \SI{5}{\second} following the final \SI{500}{\milli\second} US delivered at the end of the testing phase (health/viability check).
	\item It displayed tail movements in both the 15-s pre-CS baseline and the corresponding CR interval in at least 3 trials in each of the following trial blocks: the last 5 trials of pre-training (``late pre-train''), the first 5 trials of testing (``early test''), and the last 5 trials of testing (``late test'').
\end{itemize}

\paragraph{Number of animals included in the analysis}

A total of 87 larval zebrafish were excluded from the refined-conditioning behavioral analyses based on the predefined criteria described above. The final numbers of animals included were as follows. For the test groups: 36 fish in Delay, 34 in, 79 in 3sTrace, and 11 in 10sTrace. For the control groups: 31 fish in Delay, 36 in, 83 in 3sTrace, and 15 in 10sTrace. The corresponding numbers of excluded animals were: 12 in Delay test, 14 in test, 10 in 3sTrace test, and 6 in 10sTrace test; and 15 in Delay control, 12 in control, 8 in 3sTrace control, and 4 in 10sTrace control.

In the pharmacological experiments with the NMDA receptor antagonist MK-801, larvae were analyzed under three concentrations (0, 10, and \SI{100}{\micro\Molar}) in Delay conditioning and under two concentrations (0 and \SI{100}{\micro\Molar}) in Trace conditioning. In the Delay experiment, 14 control and 16 Delay fish were analyzed in the 0~\si{\micro\Molar} MK-801 condition, 15 control and 15 Delay fish in the 10~\si{\micro\Molar} condition, and 13 control and 13 Delay fish in the 100~\si{\micro\Molar} condition. In the Trace experiment, 13 control and 17 Trace fish were analyzed in the 0~\si{\micro\Molar} MK-801 condition and 16 control and 17 Trace fish in the 100~\si{\micro\Molar} condition. Based on the same predefined behavioral inclusion criteria, 3 control and 1 Delay fish were excluded from the 0~\si{\micro\Molar} Delay groups, 1 control and 1 Delay fish from the 10~\si{\micro\Molar} Delay groups, and 2 control and 2 Delay fish from the 100~\si{\micro\Molar} Delay groups; in the Trace experiment, 5 control and 1 Trace fish were excluded from the 0~\si{\micro\Molar} groups and 2 control and 1 Trace fish from the 100~\si{\micro\Molar} groups.

In the long-term memory experiment with the protein synthesis inhibitor puromycin, \SIrange{5}{7}{\dpf} larvae were tested either in control conditions (0~\si{\milli\gram\per\liter} puromycin) or in the presence of \SI{5}{\milli\gram\per\liter} puromycin, as in previous studies of associative LTM in larval zebrafish \citep{hinzProteinSynthesisDependentAssociative2013,dempseyRegionalSynapseGain2022}.
% TODO Add this information for the ca8 ablation experiment and adapt the following paragraph accordingly.

% TODO Add this information for the imaging experiments.

\subsection{Statistical analysis}

Statistical analyses were performed using non-parametric tests. For between-group comparisons (e.g. conditioned vs.\ control or Ntr\textsuperscript{+} vs.\ Ntr\textsuperscript{-}), we used two-sided Mann--Whitney U tests. For within-group comparisons across conditions or phases (e.g. late pre-train vs.\ early test vs.\ late test within the same group), we used Wilcoxon signed-rank tests. When multiple comparisons were performed within a family of related tests, $p$-values were adjusted using the Holm--Bonferroni correction, unless stated otherwise.

Significance annotations in figures follow the convention: ns, $p > 0.05$; *, $p \leq 0.05$; **, $p \leq 0.01$; ***, $p \leq 0.001$; ****, $p \leq 0.0001$. Where indicated, 95\% confidence intervals (CIs) for median traces or summary statistics were estimated by bootstrap resampling with 1{,}000 iterations using a fixed random seed (seed = 10) to ensure reproducibility.










% !!!! CHECK NOTES:
ONENOTE
RANDOM WORD DOCUMENTS



% endregion

% * References
% region References
% Ensure bibliography build stability without in-text citations
\bibliography{main}

% Example citation to ensure at least one \cite command is present
% Remove or replace with your actual citation as needed
% \cite{aizenbergCerebellarDependentLearningLarval2011}

% endregion


\end{document}
